\documentclass[11pt]{article}

\usepackage[a4paper,margin=1in]{geometry}
\usepackage{amsmath,amssymb,amsthm,mathtools}
\usepackage[hidelinks]{hyperref}
\hypersetup{
  pdfauthor={Bangyen Pham},
  pdftitle={The Placement Game Along (r,3,5,7): Exact Infima and Log-Periodic Oscillation},
  pdfsubject={Exact infima of a coefficient order statistic of polynomial multiples along a one-parameter family of roots},
  pdfkeywords={polynomial multiples, coefficient order statistics, reduced height, exponential sums, log-periodic oscillation},
  pdflang={en-US}
}

\newtheorem{theorem}{Theorem}[section]
\newtheorem{lemma}[theorem]{Lemma}
\newtheorem{corollary}[theorem]{Corollary}
\newtheorem{proposition}[theorem]{Proposition}
\newcommand{\R}{\mathbb R}
\newcommand{\N}{\mathbb N}
\newcommand{\E}{\mathbb E}
\DeclareMathOperator{\Tail}{Tail}
\DeclareMathOperator{\sgn}{sgn}

\title{The Placement Game Along $(r,3,5,7)$: Exact Infima and Log-Periodic Oscillation}
\author{\shortstack{Bangyen Pham\\
\small Independent Researcher\\
\small\href{mailto:bangyenp@gmail.com}{\texttt{bangyenp@gmail.com}}\\
\small\href{https://orcid.org/0009-0006-9928-2088}{ORCID: 0009-0006-9928-2088}}}
\date{September 2026}

\begin{document}
\maketitle

\begin{abstract}
For $P_r=(x-r)(x-3)(x-5)(x-7)$ we study the infimum of the second largest
nonleading coefficient magnitude over monic multiples of $P_r$, as the
root $r$ moves in $(1,3)$.  At roots at least $2$ it equals
$\tfrac{1704}{31}$, the value when the exempt coefficient escapes to the
bottom.  We show that this stays true exactly for $r\ge r_c=1.4658\ldots$,
determine the infimum in closed form on $[1.0745\ldots,r_c]$, and show that
as $r\to1^+$ it tends to $24$ and exceeds it by a quantity of exact order
$(r-1)^\alpha$, $\alpha=\log(5/3)/\log3$, whose ratio to $(r-1)^\alpha$
oscillates log-periodically between $80\cdot18^{-\alpha}$ and
$\tfrac{400}3\cdot18^{-\alpha}$.  Every cost along the family is explicit,
with a limit profile that describes the infimum uniformly near $r=1$,
across the points where the worst position of the exempt coefficient
jumps.
\end{abstract}

\noindent\textbf{Keywords.} Polynomial multiples; coefficient order
statistics; reduced height; exponential sums; log-periodic oscillation.

\noindent\textbf{2020 Mathematics Subject Classification.} Primary 11C08; Secondary 26C05, 41A50.

\section{Introduction}\label{sec:intro}

For a nonzero real polynomial $F$ let $b_k(F)$ be the $k$-th largest
magnitude of a nonleading coefficient of $F$.  For real roots
$2\le r_1\le\cdots\le r_L$ the companion paper~\cite{coefficient-mass-attainment}
computes the infimum of $b_{u+1}(QP)$ over monic $Q$, where
$P=\prod_i(x-r_i)$, and shows that below the root $2$ that computation can
fail.  This paper follows one family through the failure, $u=1$ and
\[
  P_r=(x-r)(x-3)(x-5)(x-7),\qquad 1<r<3,
\]
where every quantity turns out to be explicit.

\paragraph{Setup.}  We use the notation
of~\cite{coefficient-mass-attainment}.  Put $y_i=1/r_i$.  A
\emph{certificate} for $X\subset\N_{>0}$ on some of the nodes $y_i$ is an
exponential sum $v_d=\sum_i\lambda_iy_i^d$ with $v_0=1$ and $v|_X=0$, and
$\Tail(v)=\sum_{d\ge1}|v_d|$.  A certificate on $c$ nodes is \emph{full} if
it vanishes at $c-1$ positive integers, its \emph{zero set}; it has no other
real zeros (the \emph{zero bound}), so it keeps one sign past its largest
zero.  The \emph{cost} $\tau(X)$ of a placement $X$ is the least tail of a
certificate for $X$ on all the nodes, and $\tau^*$ is the least tail of a
sum $v$ on $y_2,\ldots,y_L$ with $v_0=1$, the cost of the \emph{escaping
placement}.  For $u=1$ the infimum is $1/T$ with
$T=\sup_x\tau(\{x\})\ge\tau^*$ \cite[Theorem~3.1]{coefficient-mass-attainment},
and $T=\tau^*$ at roots at least $2$
\cite[Theorem~4.3]{coefficient-mass-attainment}.  For
$S_j(R)$, the sum of the $j+1$ leading coefficients of a monic $R$, every
monic $M$ and finite $X$ give
$\tau(X)\ge1/\max_{j\ge1,\,j\notin X}|S_j(PM)|$
\cite[Lemma~5.1]{coefficient-mass-attainment}.  Below $2$ the reduction we
use is \cite[Theorem~5.3]{coefficient-mass-attainment}: if $V$ is a full
certificate on $y_2,\ldots,y_L$ achieving $\tau^*$, with zero set $Z^*$ and
$g=\max(Z^*\cup\{0\})$, and $H$ is the sum on all the nodes with weight $1$
on $y_1$, $H_0=0$ and $H|_{Z^*}=0$, then
\begin{enumerate}
\item[(a)] $\tau(\{x\})\le\tau^*$ for $x\in Z^*$;
\item[(b)] for $x>g$ the full certificate $w$ with zero set $Z^*\cup\{x\}$,
the \emph{aligned lift}, has
$\Tail(w)=\tau^*-2\sum_{d>x}|V_d|
-\frac{|V_x|}{|H_x|}\bigl(\sum_{d\le x}|H_d|-\sum_{d>x}|H_d|\bigr)$;
\item[(c)] with $N$ the least integer $N>g$ with
$2\sum_{d\le N}|H_d|\ge\Tail(H)$, $\tau(\{x\})<\tau^*$ for $x\ge N$ and
$T=\max\bigl(\tau^*,\max\{\tau(\{x\}):x<N,\ x\notin Z^*\}\bigr)$.
\end{enumerate}
Along $(r,3,5,7)$, $Z^*=\{1,3\}$ and
$\tau^*=\tfrac{31}{1704}$ for every $r$.

\paragraph{Results.}  The escaping placement is the worst one exactly when
$930r^2-7r-1988\ge0$, that is, $r\ge r_c=1.4658\ldots$; for
$\tfrac{13}{10}\le r\le r_c$ the infimum is $48(r-1)(15r+71)/(49r-34)$
(Theorem~\ref{thm:family}).  Below $\tfrac{13}{10}$ the zero set minimising the
cost of the worst placement keeps moving, and the worst placement changes
from $\{4\}$ to $\{5\}$; on $[1.0745\ldots,\tfrac{13}{10}]$ the infimum is
a piecewise rational function of $r$ with seven pieces, each proved by a
multiplier and explicit certificates (Theorem~\ref{thm:familylow}).  As
$r\to1^+$ the pieces accumulate: the infimum tends to $24$, the value at
$(3,5,7)$, and exceeds it by a quantity of exact order $(r-1)^\alpha$,
$\alpha=\log(5/3)/\log3$, while the worst placement is
$\log_3\frac1{r-1}+O(1)$ (Theorem~\ref{thm:familyone}); the costs that
matter are those at the roots $3,5,7$, in closed form
(Lemma~\ref{lem:threenode}), and a witness with one sign flip far out
transfers them to $r$ near $1$.

The ratio of the excess to $(r-1)^\alpha$
has no limit: off neighbourhoods of the points where $3^x(r-1)=18$ for an
integer $x$, the worst placement is the largest $x^*$ with
$3^{x^*}(r-1)<18$, its minimising zero set is $\{1,x^*,z\}$ for some
$z\ge x^*+3$, and the excess is $(80+o(1))(\tfrac35)^{x^*}$, so the ratio
oscillates log-periodically between $80\cdot18^{-\alpha}$ and
$\tfrac{400}3\cdot18^{-\alpha}$ (Theorem~\ref{thm:familyosc}), on an
explicit range of $r$ (Proposition~\ref{prop:familyrdelta}).

Every cost along the family is explicit: for each placement $\{x\}$, $x\ge2$,
and each $1<r<3$ the minimising zero set is $\{1,x,z\}$ or $\{1,z,x\}$,
with $z$ a half-mass point of an explicit exponential sum
(Theorem~\ref{thm:familyzero}).  As $3^x(r-1)\to\kappa$ the cost of
$\{x\}$ tends to an explicit function of $\kappa$, equal to
$\tfrac1{24}$ up to $\kappa=18$ and beyond it the least of the limits
$t+\ell/\kappa$ of finitely many certificates at the roots $3,5,7$ lifted
to $r$, with breakpoints accumulating at $18$
(Theorem~\ref{thm:familylimit}).  Across a jump this profile decides
between $x^*$ and $x^*+1$: as $r$ decreases to a jump the ratio falls by
the factor $\tfrac35$ within a window of width of order
$(r-1)^{\alpha/(1+\alpha)}$ ending at the jump, and off that window it
converges uniformly to the log-periodic profile
(Theorem~\ref{thm:familyjump}).  Some of these proofs end in a finite exact
computation, over $\mathbb Q$ and over $\mathbb Q(r)$, described in the
paragraph on computer assistance.

\paragraph{Organisation.}  Section~\ref{sec:familyfar} settles the family
down to $r=1.0745\ldots$, Section~\ref{sec:nearone} treats $r\to1^+$ and
the oscillation, Section~\ref{sec:onezero} computes every cost and its limit
profile, and Section~\ref{sec:jumps} crosses the jumps.
Section~\ref{sec:familyopen} lists what is open, and
Appendix~\ref{sec:familydata} the data of the rational certificates.

\section{The family away from \texorpdfstring{$r=1$}{r=1}}\label{sec:familyfar}

\begin{theorem}[The family $(r,3,5,7)$]\label{thm:family}
Let $1<r<3$, $P_r=(x-r)(x-3)(x-5)(x-7)$ and $u=1$, so $\tau^*=\tfrac{31}{1704}$,
and let $r_c=(7+\sqrt{7395409})/1860=1.46583\ldots$ be the positive root of
$930r^2-7r-1988$.  Then $T>\tau^*$ for $1<r<r_c$, and $T=\tau^*$ for
$r_c\le r<3$.  More precisely, with $\beta(r)=48(r-1)(15r+71)/(49r-34)$,
\[
  \inf_Qb_2(QP_r)\ \begin{cases}
    \le\tfrac{209}{4}, & 1<r\le\tfrac{13}{10},\\[2pt]
    =\beta(r), & \tfrac{13}{10}\le r\le r_c,\\[2pt]
    =\tfrac{1704}{31}, & r_c\le r<3,
  \end{cases}
\]
over monic $Q$, where $\beta(r)<\tfrac{1704}{31}$ for $r<r_c$.  For
$\tfrac{13}{10}\le r<r_c$, moreover, $T=\tau(\{4\})=1/\beta(r)$ and
$\tau(\{x\})<T$ for every positive integer $x\ne4$; for
$\tfrac{29}{20}\le r<r_c$ even $\tau(\{x\})\le\tau^*$.
\end{theorem}

\begin{proof}
The surviving roots are $3,5,7$ for every $r$, so $\tau^*=\tfrac{31}{1704}$,
achieved by the full certificate on $(\tfrac13,\tfrac15,\tfrac17)$ with zero
set $Z^*=\{1,3\}$ (the example after \cite[Theorem~4.3]{coefficient-mass-attainment}).  Write
$y=1/r$ for the new node.

\emph{Lower bounds.}  For $M=x+\kappa$ the partial sums of $P_rM$ are
$S_1=\kappa-r-14$, $S_2=14r+57-\kappa(r+14)$, $S_3=\kappa(14r+57)-57r-48$,
$S_j=(P_rM)(1)=48(\kappa+1)(r-1)$ for $j\ge5$, and $S_4$, which we exempt.
\cite[Lemma~5.1]{coefficient-mass-attainment} with $X=\{4\}$ gives
$T\ge\tau(\{4\})\ge1/\max_{j\ne4}|S_j|$.  At $\kappa=\tfrac32$ each $S_j$ is
affine in $r$, so $|S_j|$ is convex and its maximum on $[1,\tfrac{13}{10}]$
is taken at an end; the values are $-\tfrac{27}2,\tfrac{97}2,\tfrac32,0$ at
$r=1$ and $-\tfrac{69}5,\tfrac{209}4,-\tfrac{93}{10},36$ at
$r=\tfrac{13}{10}$, all at most $\tfrac{209}4<\tfrac{1704}{31}$.  At
$\kappa=(105-34r)/(49r-34)$,
\[
  S_2=S_5=\beta(r),\qquad
  S_1=-\frac{7(7r^2+98r-83)}{49r-34},\qquad
  S_3=-\frac{3269r^2+882r-7617}{49r-34},
\]
where $\beta(r)>0$ for $r>1$.  Since $49r-34>0$ and $7r^2+98r-83>0$ for $r\ge1$ (the
latter is $22$ at $r=1$ and increasing), $S_1<0$, so $|S_1|\le\beta(r)$
amounts to $671r^2+2002r-2827\ge0$, while $|S_3|\le\beta(r)$ amounts to
$3989r^2+3570r-11025\ge0$ and $2549r^2-1806r-4209\le0$; the first two
increase on $r>0$ and the third is convex, and all three hold at
$r=\tfrac{13}{10}$ and $r=\tfrac32$, hence on $[\tfrac{13}{10},\tfrac32]$.
There $\tau(\{4\})\ge1/\beta(r)$, and
\[
  \frac1{\beta(r)}-\tau^*=\frac{1988+7r-930r^2}{3408(r-1)(15r+71)}
\]
has the sign of $1988+7r-930r^2$, positive exactly for $r<r_c$.

\emph{Upper bounds.}  For $1<r<3$, $y>\tfrac13$ is the largest node.
Solving the $4\times4$ systems, with the sign patterns given by the zero
bound (Appendix~\ref{sec:familydata} lists the solutions), the full
certificates with zero sets $\{1,3,4\}$ and $\{1,2,5\}$ have tails
\[
  \frac1{\beta(r)}\quad\text{and}\quad
  t_2(r)=\frac{4267r^2+5070r-5871}{48(r-1)(3466r^2+7455r+11025)},
\]
and $\tau^*-t_2(r)=c_2(r)/\bigl(3408(r-1)(3466r^2+7455r+11025)\bigr)$ with
$c_2(r)=214892r^3-55639r^2-138630r-266709$, so the second is at most
$\tau^*$ exactly when $c_2(r)\ge0$.  For the $H$ of
\cite[Theorem~5.3]{coefficient-mass-attainment}, $H_d>0$ for $d\ge4$ and $H_2<0$ by the zero
bound, and solving the $3\times3$ system for $H$
(Appendix~\ref{sec:familydata}) gives, on $(1,3)$,
\[
  2\sum_{d\le5}|H_d|-\Tail(H)=-H_2+H_4+H_5-\sum_{d\ge6}H_d
  =\frac{c_5(r)(3-r)(5-r)(7-r)}{12524400\,r^5(r-1)},
\]
with $c_5(r)=206070r^3-67123r^2-161312r-238560$, which has the sign of
$c_5(r)$ there.  The two tails and this identity are identities between
rational functions of $r$, verified by exact polynomial arithmetic over
$\mathbb Q$: with $r$ an indeterminate the computation
solves the three linear systems over $\mathbb Q(r)$, forms each tail and
the left side above from the sign pattern of the zero bound and the
geometric sums past the largest zero, and compares reduced numerators and
denominators with the displayed forms, so the check covers every $r$ in
$(1,3)$, not sample points; Appendix~\ref{sec:familydata} lists the
solutions and the values entering each tail, from which the identities can
also be checked by hand.  The same test checks the partial sums and the
three quadratics of the lower bounds as identities in $r$.  Both cubics
increase on $r\ge1$, since there $r\le r^2$ and $1\le r^2$ give
$c_2'(r)=644676r^2-111278r-138630\ge394768r^2>0$ and
$c_5'(r)=618210r^2-134246r-161312\ge322652r^2>0$, and both are positive
at $\tfrac{29}{20}$, so for $\tfrac{29}{20}\le r<3$ the certificate
$\{1,2,5\}$ gives $\tau(\{2\})\le\tau^*$ and $N\le5$.  By \cite[Theorem~5.3]{coefficient-mass-attainment},
$T=\max(\tau^*,\tau(\{2\}),\tau(\{4\}))\le\max(\tau^*,1/\beta(r))$.

Now let $\tfrac{13}{10}\le r<\tfrac{29}{20}$, where $c_2$ and $c_5$ are
negative at the left end, and take one more term of $H$.  By the same sign
pattern and Appendix~\ref{sec:familydata},
\[
  2\sum_{d\le6}|H_d|-\Tail(H)=-H_2+H_4+H_5+H_6-\sum_{d\ge7}H_d
  =\frac{c_6(r)(3-r)(5-r)(7-r)}{1315062000\,r^6(r-1)}
\]
on $(1,3)$, with
$c_6(r)=23062470r^4-598283r^3-7874752r^2-16937760r-25048800$.  On $r\ge1$,
$c_6'(r)\ge(92249880-1794849-15749504-16937760)r^3=57767767r^3>0$, and
$c_6(\tfrac{13}{10})=\tfrac{522259242}{125}>0$, so $N\le6$.  The full
certificate with zero set $\{1,3,5\}$ has tail
\[
  t_5(r)=\frac{5(347r^2+250r-501)}{24(r-1)q_5(r)},\qquad
  q_5(r)=960r^2+5041r+7455,
\]
and, with $q(r)=3466r^2+7455r+11025$,
\[
  \frac1{\beta(r)}-t_2(r)=\frac{d_2(r)}{48(r-1)(15r+71)q(r)},\qquad
  \frac1{\beta(r)}-t_5(r)=\frac{d_5(r)}{16(r-1)(15r+71)q_5(r)},
\]
where $d_2(r)=105829r^3-131556r^2+14850r+41991$ and
$d_5(r)=-1670r^3-23167r^2+30517r+34080$.  For $r\ge1$,
$d_2'(r)=317487r^2-263112r+14850\ge54375r^2+14850>0$ and $d_2(1)=31114$, while
$d_5'(r)=-5010r^2-46334r+30517\le-20827$ and
$d_5(\tfrac32)=\tfrac{44187}2$; so $t_2(r)<1/\beta(r)$ and
$t_5(r)<1/\beta(r)$ on $[1,\tfrac32]$.  These three identities are checked
over $\mathbb Q(r)$ as above.  By
\cite[Theorem~5.3]{coefficient-mass-attainment}, as $\{1,3\}=Z^*$,
$T=\max(\tau^*,\tau(\{2\}),\tau(\{4\}),\tau(\{5\}))
\le\max(\tau^*,t_2(r),1/\beta(r),t_5(r))=1/\beta(r)$, since $\tau^*<1/\beta(r)$
for $r<r_c$.

For $r_c\le r<3$ the bound $T\le\max(\tau^*,1/\beta(r))$ for
$r\ge\tfrac{29}{20}$ is $\tau^*$, so $T=\tau^*$.  For
$\tfrac{29}{20}\le r<r_c$ the lower bound gives $T\ge\tau(\{4\})\ge1/\beta(r)>\tau^*$,
so $T=\tau(\{4\})=1/\beta(r)$, and every other placement costs at most
$\tau^*$: $\{1\}$ and $\{3\}$ by \cite[Theorem~5.3(a)]{coefficient-mass-attainment}, $\{2\}$ by
the certificate $\{1,2,5\}$, and $\{x\}$ with $x\ge5\ge N$ by
\cite[Theorem~5.3(c)]{coefficient-mass-attainment}.  For $\tfrac{13}{10}\le r<\tfrac{29}{20}$
likewise $T\ge\tau(\{4\})\ge1/\beta(r)$, so $T=\tau(\{4\})=1/\beta(r)$, and
every other placement costs less: $\{1\}$ and $\{3\}$ at most
$\tau^*<1/\beta(r)$, $\{2\}$ and $\{5\}$ at most $t_2(r)$ and $t_5(r)$, and
$\{x\}$ with $x\ge6\ge N$ less than $\tau^*$.  \cite[Theorem~3.1]{coefficient-mass-attainment} turns
these into the stated infima.
\end{proof}

Below $\tfrac{13}{10}$ the zero set minimising $\tau(\{4\})$ moves, and the
lower bound needs one multiplier for each zero set.  They are the
multipliers of \cite[Proposition~3.2]{coefficient-mass-attainment}, for one placement.  Let $v$
be a full certificate on $y_1,\ldots,y_L$ whose zero set $Z$ contains $x$.
Put $s_d=\sgn(v_d)$ for $d\ge1$ outside $Z$, and let $(s_d)_{d\in Z}$ and
$\theta$ solve the $L$ linear equations $\sum_{d\ge1}s_dy_i^d=\theta$,
$1\le i\le L$, where the sum past $\max Z$ is geometric; the solution is
unique, the matrix $(y_i^d)_{i\le L,\,d\in\{0\}\cup Z}$ being nonsingular by
the zero bound.  Pairing with $v$
gives $\theta=\sum_i\lambda_i\sum_{d\ge1}s_dy_i^d=\sum_{d\ge1}s_dv_d=\Tail(v)$,
since $v$ vanishes on $Z$.  The sequence $w_0=1$, $w_d=-s_d/\theta$ is
bounded and its series vanishes at every $y_i$, which is all the proof of
\cite[Lemma~5.1]{coefficient-mass-attainment} uses (and, as at the end of the proof of
\cite[Proposition~3.2]{coefficient-mass-attainment}, $w_d=S_d(PM)$ for a monic $M$).  So
\begin{equation}\label{eq:witness}
  \tau(\{x\})=\Tail(v)\qquad\text{whenever }|s_d|\le1\text{ for every }
  d\in Z\setminus\{x\}.
\end{equation}
For the placement $4$ with $Z=\{1,3,4\}$ along $(r,3,5,7)$ this is the
linear multiplier of the lower bound in the proof of
Theorem~\ref{thm:family}.

\begin{theorem}[The family below $\tfrac{13}{10}$]\label{thm:familylow}
Let $P_r=(x-r)(x-3)(x-5)(x-7)$ and $u=1$, and let $\beta$ and $t_5$ be as
in Theorem~\ref{thm:family} and its proof.  For $z=5,6,7,8$ let
$\theta_z(r)$ be the tail of the full certificate with zero set $\{1,4,z\}$
on $(\tfrac1r,\tfrac13,\tfrac15,\tfrac17)$, and put $\theta_3=1/\beta$, the
tail for $\{1,3,4\}$ (Appendix~\ref{sec:familydata} lists them).  Let
$\rho_1<\rho_2<\rho'<\rho_3<\cdots<\rho_7$ be the zeros in
$(1,\tfrac32)$ of the polynomials
$\gamma_1,\gamma_2,\gamma',\gamma_3,\ldots,\gamma_7$ of
Appendix~\ref{sec:familydata},
\begin{gather*}
  \rho_1=1.07451\ldots,\quad \rho_2=1.09547\ldots,\quad
  \rho'=1.11142\ldots,\quad \rho_3=1.12444\ldots,\quad
  \rho_4=1.13392\ldots,\\
  \rho_5=1.16696\ldots,\quad \rho_6=1.21576\ldots,\quad
  \rho_7=\tfrac{105\sqrt{4278}-1785}{3989}=1.27417\ldots .
\end{gather*}
Then for $\rho_1\le r\le\tfrac{13}{10}$
\[
  \inf_Qb_2(QP_r)=\frac1{T(r)},\qquad
  T(r)=\begin{cases}
    \theta_5(r), & \rho_1\le r\le\rho_2,\\
    t_5(r), & \rho_2\le r\le\rho_3,\\
    \theta_8(r), & \rho_3\le r\le\rho_4,\\
    \theta_7(r), & \rho_4\le r\le\rho_5,\\
    \theta_6(r), & \rho_5\le r\le\rho_6,\\
    \theta_5(r), & \rho_6\le r\le\rho_7,\\
    1/\beta(r), & \rho_7\le r\le\tfrac{13}{10}.
  \end{cases}
\]
More precisely, on $[\rho_1,\tfrac{13}{10}]$ the placement $\{5\}$ costs
$\min(\theta_5(r),t_5(r))$, which is $\theta_5(r)$ on $[\rho_1,\rho_2]$ and
$t_5(r)$ on $[\rho_2,\tfrac{13}{10}]$; on $[\rho',\tfrac{13}{10}]$ the placement
$\{4\}$ costs $\theta_8$, $\theta_7$, $\theta_6$, $\theta_5$, $1/\beta$ on
$[\rho',\rho_4]$, $[\rho_4,\rho_5]$, $[\rho_5,\rho_6]$, $[\rho_6,\rho_7]$,
$[\rho_7,\tfrac{13}{10}]$; the worst placement is $\{5\}$ for $r<\rho_3$,
$\{4\}$ for $r>\rho_3$, and both at $\rho_3$; and every other placement costs
less than $T(r)$.
\end{theorem}

With Theorem~\ref{thm:family} this gives the infimum on $[\rho_1,3)$; in
particular it is $\beta(r)$ on $[\rho_7,r_c]$, it is
$1/t_5(\tfrac{11}{10})=\tfrac{3398808}{96935}$ at $r=\tfrac{11}{10}$, as in
\cite[Theorem~5.4]{coefficient-mass-attainment}, and $1/\theta_5(\tfrac54)=\tfrac{314024}{7627}=41.172\ldots$
at $r=\tfrac54$.  Consecutive pieces agree at the breakpoints: in the
notation of Appendix~\ref{sec:familydata},
$\theta_3-\theta_5=-(r+15)\gamma_7(r)/\bigl(96(r-1)(15r+71)Q_5(r)\bigr)$,
$t_5-\theta_5=-(r+15)\gamma_2(r)/\bigl(96(r-1)q_5(r)Q_5(r)\bigr)$,
$t_5-\theta_8=-\gamma_3(r)/\bigl(96(r-1)q_5(r)Q_8(r)\bigr)$, and $\theta_z-\theta_{z+1}$
for $z=5,6,7$ is a positive rational function times $-\gamma_6$, $-\gamma_5$,
$-\gamma_4$.

\begin{proof}
Write $Z^*=\{1,3\}$, as in the proof of Theorem~\ref{thm:family}.  Every
sign of a rational function of $r$ on an interval below is checked by
counting the real zeros of its numerator and denominator with Sturm
sequences over $\mathbb Q$ between rational endpoints; together with the
closed forms, this is the computer-assisted part of the proof, described
at its end.

\emph{Lower bounds.}  Apply \eqref{eq:witness} to the placement $x=4$ with
$Z=\{1,3,4\}$ and $Z=\{1,4,z\}$, $z=5,\ldots,8$, and to $x=5$ with
$Z=\{1,3,5\}$ and $Z=\{1,4,5\}$.  In all seven cases $|s_1|<1$ on
$[\tfrac{1074}{1000},\tfrac32]$.  At the remaining zero $c\in Z\setminus\{1,x\}$,
solving the $4\times4$ systems over $\mathbb Q(r)$ gives
\[
\begin{array}{cccccc}
  x & Z & D & D(1-s_c) & D(1+s_c) & |s_c|\le1\text{ on}\\[2pt]
  4 & \{1,4,5\} & 32(r-1)Q_5 & \gamma_6 & -35r\gamma_7 & [\rho_6,\rho_7]\\
  4 & \{1,4,6\} & 16(r-1)Q_6 & \gamma_5 & -105r\gamma_6 & [\rho_5,\rho_6]\\
  4 & \{1,4,7\} & 16(r-1)Q_7 & \gamma_4 & -105r\gamma_5 & [\rho_4,\rho_5]\\
  4 & \{1,4,8\} & 32(r-1)Q_8 & \gamma' & -105r\gamma_4 & [\rho',\rho_4]\\
  4 & \{1,3,4\} & 48(r-1)(15r+71) & 4209+1806r-2549r^2 & \gamma_7 & [\rho_7,\tfrac32]\\
  5 & \{1,3,5\} & 48(r-1)q_5 & -\gamma_+ & \gamma_2 & [\rho_2,\rho_+]\\
  5 & \{1,4,5\} & 96(r-1)Q_5 & -\gamma_2 & \gamma_1 & [\rho_1,\rho_2]
\end{array}
\]
with positive denominators $D$ on $r>1$.  Each $\gamma$ has exactly one
zero in $(1,\tfrac32)$ and is negative before it and positive after it, $\rho_+=1.4831\ldots$ being the zero of
$\gamma_+$, and $4209+1806r-2549r^2$ is concave and positive at $1$ and
$\tfrac32$.  So on $(1,\tfrac32]$, $|s_c|\le1$ exactly on the interval in the
last column, and
\eqref{eq:witness} gives $\tau(\{4\})$ on $[\rho',\tfrac{13}{10}]$ and
$\tau(\{5\})=\theta_5$ on $[\rho_1,\rho_2]$, $\tau(\{5\})=t_5$ on
$[\rho_2,\tfrac{13}{10}]$, as stated; as both certificates are certificates
for $\{5\}$, $\tau(\{5\})=\min(\theta_5,t_5)$ on $[\rho_1,\tfrac{13}{10}]$.

\emph{The placements $4$ and $5$.}  By the identity for $t_5-\theta_8$
above, $\theta_8<t_5$ for $1<r<\rho_3$ and $t_5<\theta_8$ for
$\rho_3<r<\tfrac32$; moreover $\theta_8<\theta_5$ on
$[\tfrac{1074}{1000},\tfrac{118}{100}]$, and $t_5<\theta_z$ on
$[\tfrac{11242}{10000},\tfrac{13}{10}]$ for $z=3,5,6,7$, where
$\tfrac{11242}{10000}<\rho_3$.  So $\tau(\{4\})\le\theta_8<\tau(\{5\})$ for
$\rho_1\le r<\rho_3$, and $\tau(\{5\})=t_5<\tau(\{4\})$ for
$\rho_3<r\le\tfrac{13}{10}$, where $\tau(\{4\})$ is one of
$\theta_8,\theta_7,\theta_6,\theta_5,\theta_3$.

\emph{The other placements.}  Apply \cite[Theorem~5.3]{coefficient-mass-attainment} with the
escaping minimiser, $Z^*=\{1,3\}$.  With the signs $H_2<0<H_d$ ($d\ge4$) of
the proof of Theorem~\ref{thm:family},
$2\sum_{d\le n}|H_d|-\Tail(H)=-H_2+\sum_{d=4}^nH_d-\sum_{d>n}H_d$ is a
rational function of $r$, positive on $[\tfrac{1074}{1000},\tfrac{13}{10}]$
for $n=13$ and on $[b_n,\tfrac{13}{10}]$ for $n=6,\ldots,12$, where
$(b_6,\ldots,b_{12})=(1.274,1.197,1.153,1.125,1.106,1.092,1.081)$.  So
$N\le13$, and $N\le x$ once $r\ge b_x$.  The placements $\{1\}$ and $\{3\}$
cost at most $\tau^*$, and $\tau^*<\min(t_5,\theta_5)=\tau(\{5\})$ on
$[\tfrac{1074}{1000},\tfrac{13}{10}]$.  For $6\le x\le12$ and $r<b_x$ the
aligned lift, with zero set $\{1,3,x\}$, has tail below both $t_5$ and
$\theta_5$ on $[\tfrac{1074}{1000},b_x]$, so below $\tau(\{5\})$.  The
placement $\{2\}$ costs at most the tail of the full certificate
$\{1,2,13\}$, below $t_5$ and $\theta_5$ on
$[\tfrac{1074}{1000},\tfrac{1112}{1000}]$, and at most that of $\{1,2,10\}$,
below each of $\theta_3,\theta_5,\ldots,\theta_8$ on
$[\tfrac{1112}{1000},\tfrac{13}{10}]$, hence below $\tau(\{4\})$ there,
since $\rho'<\tfrac{1112}{1000}$.  So every placement other than $\{4\}$ and
$\{5\}$ costs less than $\max(\tau(\{4\}),\tau(\{5\}))$, and by
\cite[Theorem~5.3(c)]{coefficient-mass-attainment} $T$ is that maximum, which the two previous
steps evaluate as displayed.  \cite[Theorem~3.1]{coefficient-mass-attainment} gives the infimum.

The exact computation over $\mathbb Q(r)$ solves the linear systems for the
certificates and for the witnesses with $r$ an indeterminate, checks the
closed forms of the $\theta_z$, the identities displayed above and in
Appendix~\ref{sec:familydata}, and that the series of each witness vanishes
at the four nodes; it brackets each zero $\rho$ between two rationals and
checks that each $\gamma$ has exactly one zero in $(1,\tfrac32)$; and it
checks every sign on an interval quoted in this proof by a Sturm count over
$\mathbb Q$.  The tails of the certificates with zero sets $\{1,3,x\}$ and
$\{1,2,z\}$ enter only through these signs.
\end{proof}

The pieces follow a pattern.  As $r$ decreases from $r_c$ the zero set
minimising $\tau(\{4\})$ moves up, $\{1,3,4\},\{1,4,5\},\ldots,\{1,4,8\}$.
For $\{1,4,z\}$ the entry $s_z$ of the witness passes from $-1$ at the
upper end of its piece to $1$ at the lower end, where, for $z\le7$,
$\{1,4,z+1\}$ takes over with the same tail; at $\rho_7$ the entries $s_3$ of $\{1,3,4\}$ and
$s_5$ of $\{1,4,5\}$ are both $-1$.  The cost of $\{5\}$ overtakes that
of $\{4\}$ at $\rho_3$.
At $r=\tfrac{11}{10}$ the worst placement $\{5\}$ and its zero set
$\{1,3,5\}$ are those of \cite[Theorem~5.4]{coefficient-mass-attainment}, whose proof bounds
the other placements with other certificates.  Below $\rho_1$ we have no
closed form in $r$, but at each $r$ every cost is explicit
(Theorem~\ref{thm:familyzero}), and the zero sets above are its two
kinds, $\{1,x,z\}$ and $\{1,z,x\}$.  A floating-point search over zero sets in $\{1,\ldots,40\}$, at the
placements $x\le15$, suggests that the minimiser for $\{5\}$ continues
$\{1,5,6\},\{1,5,7\},\{1,5,8\},\ldots$, that the worst placement becomes $6$
near $r=1.03$ and $7$ near $r=1.01$, and that the half-mass point grows
($N=15$ at $r=1.06$).  The qualitative picture is proved below: infinitely
many pieces accumulate at $r=1$, the worst placement grows like
$\log_3\frac1{r-1}$, and the infimum tends to $24$
(Theorem~\ref{thm:familyone} and Corollary~\ref{cor:familyone}); off
neighbourhoods of the points where $3^x(r-1)=18$, the worst placement is
the largest $x$ with $3^x(r-1)<18$, with zero set $\{1,x,z\}$
(Theorem~\ref{thm:familyosc}), and near them it is that $x$ or $x+1$
(Theorem~\ref{thm:familyjump}).

\section{The family near \texorpdfstring{$r=1$}{r=1}}\label{sec:nearone}

As $r\to1^+$ the node $\tfrac1r$ tends
to $1$, where $\lambda y^d$ is constant and has infinite tail, so a cheap
certificate gives it little weight, and the costs approach those at the
roots $3,5,7$.  The limit $24$ is the
infimum at $(3,5,7)$ with $u=1$, where the surviving roots $5,7$ have
$\tau^*=\tfrac1{24}$ (\cite[Theorem~4.3]{coefficient-mass-attainment}, and the proof of
\cite[Proposition~5.5]{coefficient-mass-attainment}).  The placements at $(3,5,7)$ have
closed forms.  Below, $h_k$ is the complete homogeneous symmetric
polynomial of degree $k$ ($h_k=0$ for $k<0$); for distinct
$a_1,\ldots,a_k$ and $d\ge0$ the polynomial $p$ of degree below $k$ with
$p(a_i)=a_i^d$ satisfies
\begin{equation}\label{eq:interr}
  t^d-p(t)=h_{d-k}(a_1,\ldots,a_k,t)\prod_{i=1}^k(t-a_i),
\end{equation}
since the divided difference of $t^d$ at $a_1,\ldots,a_k,t$ is $h_{d-k}$.

\begin{lemma}[Placements at $(3,5,7)$]\label{lem:threenode}
For an integer $x\ge2$ put $A=3^{-x}$, $B=5^{-x}$, $C=7^{-x}$ and
$D=3A-10B+7C$.
\begin{enumerate}
\item[(a)] $D=\tfrac8{105}h_{x-2}(\tfrac13,\tfrac15,\tfrac17)>0$.  The full
certificate $U^{(x)}$ on $(\tfrac13,\tfrac15,\tfrac17)$ with zero set
$\{1,x\}$ has weights
\[
  \lambda_3=\frac{3(5B-7C)}{2D},\qquad
  \lambda_5=-\frac{5(3A-7C)}{2D},\qquad
  \lambda_7=\frac{7(3A-5B)}{2D},
\]
and its tail $t(x)$ satisfies
\[
  \frac1{24}-t(x)=\frac{5B-7C-45AB+84AC-35BC}{12D};
\]
$t(2)=\tfrac1{48}$ and $t(3)=\tfrac{31}{1704}$.
\item[(b)] Put
\[
  s_1=\frac{11A-60B+49C+210AB-280AC+70BC}{8D},\qquad
  s_x=\frac{1-36A+90B-56C}{6D},
\]
and $\sigma_1=s_1$, $\sigma_d=-1$ for $1<d<x$, $\sigma_x=s_x$, $\sigma_d=1$
for $d>x$.  Then $\sum_{d\ge1}\sigma_da^d=t(x)$ for
$a=\tfrac13,\tfrac15,\tfrac17$.
\item[(c)] For $x\ge4$: $0<s_1\le\tfrac12$, $0<s_x\le3^x/8$ and
$\tfrac1{24}-t(x)\ge\tfrac1{16}(\tfrac35)^x$.  For $x\ge6$ also
$\tfrac1{24}-t(x)\le\tfrac16(\tfrac35)^x$.
\item[(d)] The full certificate $W$ on $(\tfrac13,\tfrac15,\tfrac17)$ with
zero set $\{1,2\}$ has $W_d=h_{d-3}(\tfrac13,\tfrac15,\tfrac17)/105$ for
$d\ge3$ and tail $\tfrac1{48}$.  For $1<r<3$ and $x\ge3$ the full
certificate on $(\tfrac1r,\tfrac13,\tfrac15,\tfrac17)$ with zero set
$\{1,2,x\}$ has tail at most
\[
  \frac1{48}+\frac{315}{128}\cdot\frac{(r/3)^x}{r^2(r-1)} .
\]
\end{enumerate}
\end{lemma}

By (b), (c) and \eqref{eq:witness} at the roots $(3,5,7)$, $t(x)$ is the
cost of the placement $\{x\}$ there for $x\ge4$, and by (c) these costs
tend to $\tfrac1{24}$ from below at the rate $(\tfrac35)^x$.

\begin{proof}
Write $\mathbf a=(\tfrac13,\tfrac15,\tfrac17)$.

(a) By \eqref{eq:interr} at the nodes $\tfrac15,\tfrac17$ and
$t=\tfrac13$, the sum
$3^{-d}-\tfrac{10}3\,5^{-d}+\tfrac73\,7^{-d}$ equals
$(\tfrac13-\tfrac15)(\tfrac13-\tfrac17)h_{d-2}(\mathbf a)
=\tfrac8{315}h_{d-2}(\mathbf a)$; at $d=x$ it is $D/3$.  The weights are
checked by substituting them into $\sum_a\lambda_a=1$,
$\sum_a\lambda_aa=0$ and $\lambda_3A+\lambda_5B+\lambda_7C=0$.  By the zero
bound $U^{(x)}_d<0$ for $1<d<x$ and $U^{(x)}_d>0$ for $d>x$, so, with
$\sum_{d=2}^{x-1}a^d=(a^2-a^x)/(1-a)$,
\[
  t(x)=\sum_a\lambda_a\,\frac{a^x+a^{x+1}-a^2}{1-a},
\]
which is the stated closed form; at $x=2,3$ it gives $\tfrac1{48}$ and
$\tfrac{31}{1704}$, the latter as in the example after
\cite[Theorem~4.3]{coefficient-mass-attainment}.

(b) For $a\in(0,1)$,
$\sum_{d\ge1}\sigma_da^d=s_1a-(a^2-a^x)/(1-a)+s_xa^x+a^{x+1}/(1-a)$, and
at $a=\tfrac13,\tfrac15,\tfrac17$ the identity is checked by substitution,
$a^x$ being $A$, $B$, $C$.

(c) Put $\beta=(\tfrac35)^x=B/A$ and $\gamma=(\tfrac37)^x=C/A$, so
$\gamma=(\tfrac57)^x\beta$ and $D=A(3-10\beta+7\gamma)$.  For $x\ge4$,
$\beta\le(\tfrac35)^4$, $\gamma\le(\tfrac37)^4$, $(\tfrac57)^x\le(\tfrac57)^4$
and $A\le3^{-4}$, so $3-10\beta+7\gamma\ge3-10(\tfrac35)^4>\tfrac{17}{10}$.
The numerator of $s_x$ lies between $1-36\cdot3^{-4}-56\cdot7^{-4}>0$ and
$1+90\cdot5^{-4}$, and
$(1+90\cdot5^{-4})/(6(3-10(\tfrac35)^4))<\tfrac18$, so $0<s_x\le3^x/8$.
Dividing by $A$,
\begin{align*}
  8(3-10\beta+7\gamma)\,s_1&=11-60\beta+49\gamma+A(210\beta-280\gamma+70\beta\gamma),\\
  8(3-10\beta+7\gamma)(\tfrac12-s_1)&=1+20\beta-21\gamma-A(210\beta-280\gamma+70\beta\gamma),\\
  12(3-10\beta+7\gamma)(\tfrac1{24}-t(x))&=5\beta-7\gamma-A(45\beta-84\gamma+35\beta\gamma).
\end{align*}
The first is at least $11-60(\tfrac35)^4-280\cdot3^{-4}(\tfrac37)^4>0$, the
second at least
$\beta\bigl(20-21(\tfrac57)^4-210\cdot3^{-4}-70\cdot3^{-4}(\tfrac37)^4\bigr)>0$,
and the third at least
$\beta\bigl(5-7(\tfrac57)^4-45\cdot3^{-4}-35\cdot3^{-4}(\tfrac37)^4\bigr)$,
which with $3-10\beta+7\gamma\le3+7(\tfrac37)^4$ gives
$\tfrac1{24}-t(x)\ge\tfrac1{16}\beta$.  For $x\ge6$ the third is at most
$5\beta+84A\gamma\le\beta(5+84\cdot3^{-6}(\tfrac57)^6)$, and
$3-10\beta+7\gamma\ge3-10(\tfrac35)^6$ gives
$\tfrac1{24}-t(x)\le\tfrac16\beta$.

(d) By \eqref{eq:interr} at the nodes $\mathbf a$ and $t=0$, the value
$W_d=p(0)$, a combination of $3^{-d},5^{-d},7^{-d}$, has $W_0=1$,
$W_1=W_2=0$ and $W_d=\tfrac1{105}h_{d-3}(\mathbf a)>0$ for $d\ge3$, so
$\Tail(W)=\tfrac1{105}\sum_kh_k(\mathbf a)=\tfrac1{105}\prod_a(1-a)^{-1}=\tfrac1{48}$.
Let $y=\tfrac1r\in(\tfrac13,1)$.  By \eqref{eq:interr} at $t=y$, the
exponential sum $H'_d=y^d-p(y)$ on the four nodes, with weight $1$ on $y$,
vanishes at $d=0,1,2$ and equals
$h_{d-3}(\mathbf a,y)\prod_a(y-a)>0$ for $d\ge3$.  So
$w=W-(W_x/H'_x)H'$ is the full certificate with zero set $\{1,2,x\}$, and
$\Tail(w)\le\tfrac1{48}+(W_x/H'_x)\Tail(H')$.  Here
$h_k(\mathbf a)=3^{-k}h_k(1,\tfrac35,\tfrac37)\le3^{-k}/((1-\tfrac35)(1-\tfrac37))=\tfrac{35}8\,3^{-k}$,
so $W_x\le\tfrac1{24}3^{3-x}$; $H'_x\ge y^{x-3}\prod_a(y-a)$; and
$\Tail(H')=\prod_a(y-a)\sum_kh_k(\mathbf a,y)
=\prod_a(y-a)/\bigl((1-y)\prod_a(1-a)\bigr)$.  Hence
$(W_x/H'_x)\Tail(H')\le\tfrac1{24}3^{3-x}\cdot\tfrac{105}{48}\cdot
r^{x-3}\cdot\tfrac r{r-1}$, the stated bound.
\end{proof}

\begin{theorem}[The family near $r=1$]\label{thm:familyone}
Let $P_r=(x-r)(x-3)(x-5)(x-7)$ and $u=1$, let $\tau(\{x\})$ be the costs at
$P_r$ and $t(x)$ as in Lemma~\ref{lem:threenode}, and put
$\alpha=\log(5/3)/\log3=0.46497\ldots$.
\begin{enumerate}
\item[(a)] For $1<r<3$, $\inf_Qb_2(QP_r)>24$.
\item[(b)] For every $x\ge4$, $\tau(\{x\})\to t(x)$ as $r\to1^+$.
\item[(c)] For $1<r\le\tfrac{101}{100}$, with $\varepsilon=r-1$,
\[
  24+\tfrac32\,\varepsilon^\alpha\le\inf_Qb_2(QP_r)\le24+125\,\varepsilon^\alpha,
\]
and every worst placement $\{x\}$, that is $\tau(\{x\})=T$, satisfies
$\log_3\frac1\varepsilon-3<x<\log_3\frac1\varepsilon+6$.
\end{enumerate}
\end{theorem}

\begin{proof}
Let $y=\tfrac1r$.  The certificates $U^{(x)}$ and $W$ of
Lemma~\ref{lem:threenode}, with weight $0$ on $y$, are certificates on the
four nodes, so $\tau(\{x\})\le t(x)$ for $x\ge2$ and
$\tau(\{1\}),\tau(\{2\})\le\tfrac1{48}$.

(a) By Lemma~\ref{lem:threenode}, $\tfrac1{48}$, $t(3)=\tfrac{31}{1704}=\tau^*$
and $t(x)$ for $x\ge4$ are below $\tfrac1{24}$.  By
\cite[Theorem~5.3(c)]{coefficient-mass-attainment}, $T$ is the largest of $\tau^*$ and finitely
many $\tau(\{x\})$, so $T<\tfrac1{24}$, and \cite[Theorem~3.1]{coefficient-mass-attainment} gives
the claim.

\emph{The witness.}  We show: if $x\ge4$, $z_0\ge x+3$ is an integer, and
$1<r<3$ satisfies
\[
  \text{(W1)}\ \ y^x(1+y)-y^2\ge\tfrac{1-y}{24},\qquad
  \text{(W2)}\ \ y^2+2y^{z_0}-y^x-y^{x+1}\ge(1-y)\bigl(1+\tfrac{3^x}8\bigr),
\]
then $\tau(\{x\})\ge t(x)/(1+3^{6-z_0})$.  Let $\sigma$ be as in
Lemma~\ref{lem:threenode}(b) and $t=t(x)$.  For an integer $z\ge z_0$ and
$\varphi\in[0,1]$ let $\rho$ agree with $\sigma$ below $z$, with
$\rho_z=1-2\varphi$ and $\rho_d=-1$ for $d>z$.  As $\sigma_d=1$ for $d\ge z$,
$\sum_{d\ge1}\rho_da^d=t-\delta(a)$ at $a=\tfrac13,\tfrac15,\tfrac17$, where
$\delta(a)=2\varphi a^z+2a^{z+1}/(1-a)\in[0,2a^z/(1-a)]$.  Let
$c_1,c_2,c_3$ solve $\sum_{j\le3}c_ja^j=\delta(a)$ at the three nodes.  The
inverse of the matrix $(a^j)$, rows indexed by $j$, has rows
$(\tfrac{27}8,-\tfrac{125}4,\tfrac{343}8)$,
$(-\tfrac{81}2,\tfrac{625}2,-343)$,
$(\tfrac{945}8,-\tfrac{2625}4,\tfrac{5145}8)$, and as $z\ge7$ this gives
$|c_1|\le13\cdot3^{-z}$, $|c_2|\le146\cdot3^{-z}$, $|c_3|\le405\cdot3^{-z}$,
so $\sum_j|c_j|<\tfrac12$ and $|c_j|\le3^{6-z}$.  Put $\rho'_j=\rho_j+c_j$
for $j\le3$, positions below $x$, and $\rho'_d=\rho_d$ otherwise; then
$\sum_{d\ge1}\rho'_da^d=t$ at the three nodes.  Order the pairs
$(z,\varphi)$ by increasing $z$ and, for each $z$, decreasing $\varphi$;
$(z,0)$ and $(z+1,1)$ give the same $\rho$ and $c$, so along this path
$F=\sum_{d\ge1}\rho'_dy^d-t$ is continuous.  At $(z_0,1)$,
\[
  \sum_{d\ge1}\rho_dy^d=s_1y+s_xy^x-\frac{y^2+2y^{z_0}-y^x-y^{x+1}}{1-y},
\]
and $s_1\le\tfrac12$, $0<s_xy^x<3^x/8$, $\sum_j|c_j|<\tfrac12$ and (W2) give
$F<-t<0$.  As $z\to\infty$, $c\to0$ and
$F\to\sum_d\sigma_dy^d-t\ge(y^x+y^{x+1}-y^2)/(1-y)-t\ge\tfrac1{24}-t>0$ by
$s_1,s_x>0$ and (W1).  So $F=0$ at some $(z,\varphi)$ with $z\ge z_0$.  There
put $w_0=1$ and $w_d=-\rho'_d/t$ for $d\ge1$: $w$ is bounded and its series
$1-\tfrac1t\sum_d\rho'_da^d$ vanishes at all four nodes, which is all the
proof of \cite[Lemma~5.1]{coefficient-mass-attainment} uses.  Off the position $x$,
$|\rho'_1|\le\tfrac12+\tfrac12$, $|\rho'_2|,|\rho'_3|\le1+3^{6-z}$,
$|\rho'_z|\le1$ and $|\rho'_d|=1$ otherwise, so every certificate for
$\{x\}$ has tail at least $t/(1+3^{6-z})\ge t/(1+3^{6-z_0})$.

(b) Fix $x\ge4$ and $z_0\ge x+3$.  As $r\to1^+$ the left sides of (W1) and
(W2) tend to $1$ and the right sides to $0$, so
$t(x)/(1+3^{6-z_0})\le\tau(\{x\})\le t(x)$ for $r$ near $1$; and $z_0$ is
arbitrary.

(c) \emph{Upper bound.}  Let $x_r$ be the largest integer with
$3^{x_r}\le8/(r\varepsilon)-48$.  As $r\varepsilon\le\tfrac{101}{10000}<\tfrac8{777}$,
$x_r\ge6$; as $8/r-48\varepsilon\ge3$, $3^{x_r}>(8/(r\varepsilon)-48)/3\ge1/\varepsilon$,
so $(\tfrac35)^{x_r}=3^{-\alpha x_r}\le\varepsilon^\alpha$.  Take $x=x_r$ and
$z_0=x+3$.  The left side of (W2) is
$y^2-(1-y)y^x(1+2y+2y^2)\ge y^2-5(1-y)$, and
$y^2/(1-y)=1/(r\varepsilon)\ge6+3^x/8$, which gives (W2).  For (W1),
$x\log r\le\varepsilon\log_3(8/\varepsilon)\le\tfrac1{100}\log_3800<\tfrac1{14}$,
as $\varepsilon\log(8/\varepsilon)$ increases on $(0,\tfrac1{100}]$, so
$y^{x-2}\ge y^x>\tfrac{13}{14}$ and
$y^x(1+y)-y^2\ge y^2(\tfrac{13}{14}\cdot\tfrac{201}{101}-1)>\tfrac{1-y}{24}$.
So $\tau(\{x_r\})\ge t(x_r)/(1+3^{3-x_r})$.  By
Lemma~\ref{lem:threenode}(c), $24(\tfrac1{24}-t(x_r))\le4(\tfrac35)^{x_r}<\tfrac15$,
so $1/t(x_r)\le24+30\cdot4(\tfrac35)^{x_r}<30$, and
\[
  \inf_Qb_2(QP_r)\le\frac{1+3^{3-x_r}}{t(x_r)}
  \le24+120(\tfrac35)^{x_r}+810\cdot3^{-x_r}\le24+121(\tfrac35)^{x_r},
\]
as $810\cdot5^{-x_r}<1$; this is at most $24+125\varepsilon^\alpha$.

\emph{Lower bound.}  Let $X$ be the least integer $X\ge4$ with
$\tfrac{315}{128}(r/3)^{X+1}\le r^2\varepsilon/192$.  By
Lemma~\ref{lem:threenode}(d) every placement $x>X$ costs at most
$\tfrac1{48}+\tfrac1{192}=\tfrac5{192}$; the placements $x\le3$ cost at most
$\tfrac1{48}$ or $\tfrac{31}{1704}$, and $\tau^*=\tfrac{31}{1704}$; and
$4\le x\le X$ costs at most $t(x)\le\tfrac1{24}-\tfrac1{16}(\tfrac35)^X$.  As
$\tfrac1{16}(\tfrac35)^X\le\tfrac1{16}(\tfrac35)^4<\tfrac1{24}-\tfrac5{192}$,
$T\le\tfrac1{24}-\tfrac1{16}(\tfrac35)^X$, and
$\inf_Qb_2(QP_r)=1/T\ge24+36(\tfrac35)^X$.  If $X=4$, $(\tfrac35)^X>\tfrac1{23}$.
If $X>4$, minimality gives $(3/r)^X<\tfrac{945}{2r^2\varepsilon}\le\tfrac{945}{2\varepsilon}$,
so $X<\log\tfrac{945}{2\varepsilon}/\log\tfrac{300}{101}$; as
$\log r\le\varepsilon$, $\varepsilon\log\tfrac{945}{2\varepsilon}\le\tfrac1{100}\log47250<\tfrac{27}{250}$
(the left side increases on $(0,\tfrac1{100}]$) and $\log\tfrac{300}{101}>\tfrac{27}{25}$,
$X\log r<\tfrac1{10}$ and $r^X<e^{1/10}<\tfrac{111}{100}$, so $3^X<\tfrac{111}{100}\cdot\tfrac{945}{2\varepsilon}<\tfrac{525}\varepsilon$
and $(\tfrac35)^X=3^{-\alpha X}>(\varepsilon/525)^\alpha>\varepsilon^\alpha/23$,
as $\alpha<\tfrac12$ ($25<27$).  Either way $\inf_Qb_2(QP_r)\ge24+\tfrac{36}{23}\varepsilon^\alpha$.

\emph{The worst placement.}  As $r<r_c$, $T>\tau^*$
(Theorem~\ref{thm:family}), so $T=\tau(\{x\})$ for some $x$
(\cite[Proposition~5.2]{coefficient-mass-attainment}).  By the upper bound,
$T\ge\tau(\{x_r\})\ge(\tfrac1{24}-\tfrac16(\tfrac35)^6)/(1+\tfrac1{27})>\tfrac5{192}$,
so a worst placement has $4\le x\le X$; and $X<\log_3\tfrac{525}\varepsilon<\log_3\tfrac1\varepsilon+6$
when $X>4$, while $4<\log_3\tfrac1\varepsilon$.  Also
$T\ge1/(24+125\varepsilon^\alpha)\ge\tfrac1{24}-\tfrac{125}{576}\varepsilon^\alpha$
and $T=\tau(\{x\})\le t(x)\le\tfrac1{24}-\tfrac1{16}(\tfrac35)^x$, so
$3^{-\alpha x}=(\tfrac35)^x\le\tfrac{125}{36}\varepsilon^\alpha$ and
$x\ge\log_3\tfrac1\varepsilon-\log_3(\tfrac{125}{36})/\alpha>\log_3\tfrac1\varepsilon-3$,
as $\log_3\tfrac{125}{36}<\tfrac65$ and $\alpha>\tfrac25$
($(\tfrac53)^5>9$).
\end{proof}

The closed forms and inequalities of Lemma~\ref{lem:threenode}, the
constants $13,146,405$, the bound $(1+3^{3-x})/t(x)\le24+121(\tfrac35)^x$
for $6\le x\le60$, and the witness itself at $r=\tfrac{101}{100}$ and
$x=x_r=6$, where $F$ vanishes at $z=41$, are checked exactly.  The
witness needs $3^x$ of order at most $1/\varepsilon$: the entry $s_x$,
about $3^x/18$, must be balanced at the node $y$ by the flip, whose weight
there is at most $y^{x+1}/(1-y)$.  Beyond that the node $\tfrac1r$ makes
the placement cheap: at $r=\tfrac{101}{100}$ the certificate $\{1,2,9\}$
costs less than $t(9)/(1+3^{-6})$, which the same test checks as a control.

\begin{corollary}\label{cor:familyone}
Along $(r,3,5,7)$ with $u=1$, $\inf_Qb_2(QP_r)\to24$ as $r\to1^+$, and the
worst placement tends to $\infty$.  On no interval $(1,1+\delta)$ does
$T(r)$ agree at every $r$ with one of finitely many rational functions of
$r$.  In particular, as the tail of the full certificate with a given zero
set is a rational function of $r$ on $(1,3)$, infinitely many pieces as in
Theorem~\ref{thm:familylow} accumulate at $r=1$.
\end{corollary}

\begin{proof}
The first two claims are Theorem~\ref{thm:familyone}(c).  By it,
$\tfrac1{24}-T=(\inf_Qb_2-24)/(24\inf_Qb_2)$ lies between
$c\varepsilon^\alpha$ and $C\varepsilon^\alpha$ for some $c,C>0$ and all
$\varepsilon\le\tfrac1{100}$.  If $T$ took values among rational functions
$f_1,\ldots,f_k$ on $(1,1+\delta)$, one $f_j$ would equal $T$ at points
$r_n\to1^+$.  Then $f_j-\tfrac1{24}$ is a rational function, nonzero at the
$r_n$ and tending to $0$ along them, so it has a zero of some order $m\ge1$
at $1$ and is $O(r-1)$ there, against $\tfrac1{24}-T\ge c\varepsilon^\alpha$
with $\alpha<1$.  The tail is rational in $r$ because the weights solve a
linear system with entries $r^{-d}$, the zero bound fixes the sign pattern,
and the part past the largest zero is geometric.
\end{proof}

\paragraph{The constant near $r=1$.}  The constants $\tfrac32$ and $125$ of
Theorem~\ref{thm:familyone}(c) cannot be brought to a common value:
$(\inf_Qb_2(QP_r)-24)(r-1)^{-\alpha}$ oscillates log-periodically between
$80\cdot18^{-\alpha}$ and $\tfrac{400}3\cdot18^{-\alpha}$
(Theorem~\ref{thm:familyosc}).  The switch is at $3^x(r-1)=18$.  Below it
the flipped witness works with the true size of $s_x$, which is $3^x/18$ up
to a factor tending to $1$, and the placement $\{x\}$ costs $t(x)$ up to
$O(3^{-x})$; above it a small weight on $\tfrac1r$ added to a certificate
$U^{(c)}$ at $(3,5,7)$ makes the placement cheaper by an amount that does
not tend to $0$.  For $c\ge2$ let $\ell(c)$ and $\lambda_7(c)$ be the weights
$\lambda_3$ and $\lambda_7$ of the certificate $U^{(c)}$ of
Lemma~\ref{lem:threenode}(a), and put
\[
  K_c=\frac{\ell(c)}{\tfrac1{24}-t(c)},\qquad
  R(x)=\tfrac{36}5\bigl(\tfrac53\bigr)^x\bigl(\tfrac1{24}-t(x)\bigr).
\]

\begin{lemma}[Both sides of $3^x(r-1)=18$]\label{lem:eighteen}
Let $1<r<3$, $\varepsilon=r-1$, $y=\tfrac1r$, and let $\tau(\{x\})$ be the
costs at $P_r=(x-r)(x-3)(x-5)(x-7)$.
\begin{enumerate}
\item[(a)] Let $x\ge4$, and let $t=t(x)$, $s_1$ and $s_x$ be as in
Lemma~\ref{lem:threenode}.  If
\[
  y^x(1+y)-y^2\ge\frac{1-y}{24}\qquad\text{and}\qquad
  s_x+7\le\frac1{r\varepsilon},
\]
then for some integer $z\ge x+3$ the cost $\tau(\{x\})$ is the tail of the
full certificate with zero set $\{1,x,z\}$, and
$t(x)-14\cdot3^{-z}\le\tau(\{x\})\le t(x)$.
\item[(b)] For integers $c\ge2$ and $x>c$,
\[
  \tau(\{x\})\le t(c)+\frac{35}{24}\cdot
  \frac{\ell(c)3^{-x}+\lambda_7(c)7^{-x}}{h_{x-2}(\tfrac15,\tfrac17,1)}\cdot
  \frac{r^{x-1}}{\varepsilon},
\]
and the right side decreases as $x$ increases.
\item[(c)] As $x\to\infty$, $18\cdot3^{-x}s_x\to1$ and $R(x)\to1$; and
$R(x)\ge\tfrac9{20}$ for $x\ge4$.  For $c\ge2$, $K_c>18$, and $K_c\to18$ as
$c\to\infty$.
\end{enumerate}
\end{lemma}

\begin{proof}
Write $A=3^{-x}$, $B=5^{-x}$, $C=7^{-x}$ and $D=3A-10B+7C$, as in
Lemma~\ref{lem:threenode}.

(a) We run the witness of the proof of Theorem~\ref{thm:familyone} again,
correcting it at the positions $1$ and $x$ instead of $1,2,3$, so that no
entry off $x$ leaves $[-1,1]$.  Let $\sigma$ be as in
Lemma~\ref{lem:threenode}(b).  For an integer $z\ge x+3$ and
$\varphi\in[0,1]$ put $\delta(a)=2\varphi a^z+2a^{z+1}/(1-a)$, which lies in
$[0,2a^z/(1-a)]$, and let $p,q,\theta'$ solve
$pa+qa^x-\theta'=\delta(a)$ at $a=\tfrac13,\tfrac15,\tfrac17$.  Subtracting
the equations in pairs and solving,
\[
  q=\frac{3\delta(\tfrac13)-10\delta(\tfrac15)+7\delta(\tfrac17)}{D},\qquad
  p=-\frac{105}{2D}\Bigl(\delta(\tfrac13)(B-C)-\delta(\tfrac15)(A-C)
  +\delta(\tfrac17)(A-B)\Bigr),
\]
and $\theta'=\tfrac p3+qA-\delta(\tfrac13)$.  Here $D\ge\tfrac{17}{10}A$
(proof of Lemma~\ref{lem:threenode}(c)), $0\le B-C\le B$,
$0\le A-B,A-C\le A$ and $B\le(\tfrac35)^4A$, so, as $z\ge7$,
\[
  |p|\le\tfrac{525}{17}\,3^{-z}\bigl(3(\tfrac35)^4+\tfrac52(\tfrac35)^7
  +\tfrac73(\tfrac37)^7\bigr)\le15\cdot3^{-z};
\]
and $|q|\le\max\bigl(3\delta(\tfrac13)+7\delta(\tfrac17),10\delta(\tfrac15)\bigr)/D
\le\tfrac{10}{17}\,3^{x-z}\bigl(9+\tfrac{49}3(\tfrac37)^7\bigr)\le6\cdot3^{x-z}$,
as $25\cdot5^{-z}\le9\cdot3^{-z}$.  Hence $|\theta'|\le14\cdot3^{-z}$.  Let
$\rho$ agree with $\sigma$ except that $\rho_1=s_1+p$, $\rho_x=s_x+q$,
$\rho_z=1-2\varphi$ and $\rho_d=-1$ for $d>z$.  As $\sigma_d=1$ for $d>x$,
$\sum_{d\ge1}\rho_da^d=t-\delta(a)+pa+qa^x=t+\theta'$ at
$a=\tfrac13,\tfrac15,\tfrac17$.  Put
$F(z,\varphi)=\sum_{d\ge1}\rho_dy^d-t-\theta'$.  The pairs $(z,0)$ and
$(z+1,1)$ give the same $\delta$, hence the same $p,q,\theta'$ and $\rho$,
so $F$ is continuous along the path through the pairs ordered by increasing
$z$ and, for each $z$, decreasing $\varphi$.  At $(x+3,1)$,
\[
  \sum_{d\ge1}\rho_dy^d=(s_1+p)y+(s_x+q)y^x
  -\frac{y^2+2y^{x+3}-y^x-y^{x+1}}{1-y},
\]
where $y^2+2y^{x+3}-y^x-y^{x+1}=y^2-(1-y)y^x(2y^2+2y+1)\ge y^2-5(1-y)$ and
$y^2/(1-y)=1/(r\varepsilon)$.  With $0<s_1\le\tfrac12$ and $s_x>0$
(Lemma~\ref{lem:threenode}(c)), $|p|\le15\cdot3^{-7}$, $|q|\le6\cdot3^{-3}$
and $t+\theta'>0$, this gives $F(x+3,1)<s_x+7-1/(r\varepsilon)\le0$.  As
$z\to\infty$, $p,q,\theta'\to0$ and
$F\to\sum_d\sigma_dy^d-t\ge(y^x+y^{x+1}-y^2)/(1-y)-t\ge\tfrac1{24}-t>0$, as in
the proof of Theorem~\ref{thm:familyone}.  So $F=0$ at some $(z,\varphi)$
with $z\ge x+3$.  There $\theta=t+\theta'$ satisfies
$\sum_{d\ge1}\rho_da^d=\theta$ at all four nodes,
$|\rho_1|\le\tfrac12+15\cdot3^{-7}<1$, $|\rho_z|\le1$, and $|\rho_d|=1$ for
every other $d\ne x$.  For a certificate $v'$ for $\{x\}$, summing over the
nodes as in the proof of \cite[Lemma~5.1]{coefficient-mass-attainment} gives
$\theta=\sum_{d\ge1}\rho_dv'_d\le\Tail(v')$, so $\tau(\{x\})\ge\theta$.  The
full certificate $v$ with zero set $\{1,x,z\}$ is positive at $0$ and
changes sign exactly at $1$, $x$ and $z$, so $\rho_d=\sgn(v_d)$ off its zero
set, and the same sum gives $\theta=\Tail(v)$.  Hence
$\tau(\{x\})=\Tail(v)=\theta\ge t-14\cdot3^{-z}$, while $\tau(\{x\})\le t$
as $U^{(x)}$ is a certificate for $\{x\}$.

(b) By \eqref{eq:interr} at the nodes $\tfrac15,\tfrac17$ and $t=y$, the
sum $g_d=y^d-\tfrac{35y-5}2\,5^{-d}+\tfrac{35y-7}2\,7^{-d}$ equals
$(y-\tfrac15)(y-\tfrac17)h_{d-2}(\tfrac15,\tfrac17,y)$: it vanishes at
$d=0,1$, is positive for $d\ge2$, and
$\sum_dg_d=(y-\tfrac15)(y-\tfrac17)/\bigl((1-\tfrac15)(1-\tfrac17)(1-y)\bigr)$.
For $x>c$, $U^{(c)}-(U^{(c)}_x/g_x)g$ is a certificate for $\{x\}$ on the
four nodes, so
\[
  \tau(\{x\})\le t(c)+\frac{|U^{(c)}_x|}{g_x}\sum_dg_d
  =t(c)+\frac{35}{24}\cdot\frac{|U^{(c)}_x|}{(1-y)h_{x-2}(\tfrac15,\tfrac17,y)} .
\]
Here $h_k(\tfrac15,\tfrac17,y)=y^k\sum_{i+j\le k}(5y)^{-i}(7y)^{-j}\ge
y^kh_k(\tfrac15,\tfrac17,1)$ and $(1-y)y^{x-2}=\varepsilon r^{1-x}$;
and $0<U^{(c)}_x\le\ell(c)3^{-x}+\lambda_7(c)7^{-x}$, as $U^{(c)}$ is
positive past $c$ and its weight on $\tfrac15$ is negative by
Lemma~\ref{lem:threenode}(a).  The bound decreases in $x$, as $r<3$ and
$h_k(\tfrac15,\tfrac17,1)$ increases with $k$.

(c) Dividing the closed forms of Lemma~\ref{lem:threenode} by $A$,
\[
  18\cdot3^{-x}s_x=\frac{1-36A+90B-56C}{1-\tfrac{10}3(\tfrac35)^x+\tfrac73(\tfrac37)^x},
  \qquad
  R(x)=\frac{1-\tfrac75(\tfrac57)^x-9A+\tfrac{84}5A(\tfrac57)^x-7C}
  {1-\tfrac{10}3(\tfrac35)^x+\tfrac73(\tfrac37)^x},
\]
and both tend to $1$; $R(x)\ge\tfrac9{20}$ is
$\tfrac1{24}-t(x)\ge\tfrac1{16}(\tfrac35)^x$.  At $x=c$,
$K_c=18(5B-7C)/(5B-7C-45AB+84AC-35BC)$.  The denominator is
$12D(\tfrac1{24}-t(c))>0$ and falls short of $5B-7C$ by
$A(45B-84C)+35BC>0$, as $(\tfrac75)^c\ge\tfrac{49}{25}>\tfrac{84}{45}$; so
$K_c>18$.  And $K_c\to18$, since
$(45AB+84AC+35BC)/(5B-7C)\le(45A+84A(\tfrac57)^c+35C)/(5-7(\tfrac57)^c)\to0$.
\end{proof}

As $h_k(\tfrac15,\tfrac17,1)<\tfrac{35}{24}$, the bound in (b) exceeds
$t(c)+\ell(c)/(3^x\varepsilon)$, which exceeds $\tfrac1{24}$ while
$3^x\varepsilon<18$, as $K_c>18$: there only (a) is informative.  Once
$3^x\varepsilon\ge\kappa>18$ the bound tends, as $r\to1^+$, to at most
$t(c)+\ell(c)/\kappa$, below $\tfrac1{24}$ for $c$ with $K_c<\kappa$.

\begin{theorem}[The constant oscillates]\label{thm:familyosc}
Let $P_r=(x-r)(x-3)(x-5)(x-7)$, $u=1$, $\varepsilon=r-1$ and
$\alpha=\log(5/3)/\log3$, and put
\[
  \Delta(r)=\bigl(\inf_Qb_2(QP_r)-24\bigr)\varepsilon^{-\alpha},\qquad
  s=\log_3\frac{18}\varepsilon,\qquad x^*=\lceil s\rceil-1,\qquad
  \omega=s-x^*,
\]
so that $x^*$ is the largest integer with $3^{x^*}\varepsilon<18$ and
$0<\omega\le1$, and $\Phi(\omega)=80\cdot18^{-\alpha}(\tfrac53)^\omega$, so
that $\Phi(\omega)=80(\tfrac35)^{x^*}\varepsilon^{-\alpha}$.
\begin{enumerate}
\item[(a)] Let $0<\delta<\tfrac12$.  For $r>1$ close enough to $1$ with
$\delta\le\omega\le1-\delta$, the placement $\{x^*\}$ is the only worst
placement, $T=\tau(\{x^*\})$ is the tail of the full certificate with zero
set $\{1,x^*,z\}$ for some $z\ge x^*+3$, and
\[
  \frac1{t(x^*)}\le\inf_Qb_2(QP_r)\le\frac1{t(x^*)-14\cdot3^{-x^*-3}} .
\]
Moreover $\Delta(r)-\Phi(\omega)\to0$ as $r\to1^+$ with
$\delta\le\omega\le1-\delta$.
\item[(b)] $\liminf_{r\to1^+}\Delta(r)=\Phi(0)=80\cdot18^{-\alpha}=20.865\ldots$
and $\limsup_{r\to1^+}\Delta(r)=\Phi(1)=\tfrac{400}3\cdot18^{-\alpha}=34.775\ldots$,
and the limit points of $\Delta(r)$ as $r\to1^+$ fill $[\Phi(0),\Phi(1)]$.
\end{enumerate}
\end{theorem}

So $\Delta$ has no limit.  Along $r_n=1+18\cdot3^{-n-\omega_0}$ it tends to
$\Phi(\omega_0)$, a function of $\log_3\frac1{r-1}$ modulo $1$: the excess
$\inf_Qb_2-24$ is $(80+o(1))(\tfrac35)^{x^*}$ off the jumps, and it drops
by the factor $\tfrac35$ each time $3^x(r-1)$ crosses $18$ as $r$ decreases.

\begin{proof}
As $r\to1^+$, $\varepsilon\to0$, $s\to\infty$ and $s\varepsilon\to0$.  By
Theorem~\ref{thm:familyone} and Corollary~\ref{cor:familyone}, $T<\tfrac1{24}$,
$T\to\tfrac1{24}$ and $(\tfrac1{24}-T)\varepsilon^{-\alpha}$ is bounded, so
$\inf_Qb_2-24=24(\tfrac1{24}-T)/T$ gives
$\Delta(r)-576(\tfrac1{24}-T)\varepsilon^{-\alpha}\to0$.  As
$(\tfrac35)^x=3^{-\alpha x}$, $3^\alpha=\tfrac53$ and
$\varepsilon=18\cdot3^{-s}$,
\begin{equation}\label{eq:oscscale}
  80\,(\tfrac35)^x\varepsilon^{-\alpha}=\Phi(0)\,(\tfrac53)^{s-x}
  \qquad\text{for every }x .
\end{equation}
Recall from the proof of Theorem~\ref{thm:familyone} that
$\tau(\{1\}),\tau(\{2\})\le\tfrac1{48}$ and $\tau(\{x\})\le t(x)$ for
$x\ge3$, where $t(3)=\tau^*<\tfrac1{48}$ and
$\tfrac1{24}-t(x)=\tfrac5{36}(\tfrac35)^xR(x)$ for $x\ge4$.  Fix
$\delta\in(0,1)$.

\emph{(i) Far placements.}  There is $\eta>0$ with
$\tau(\{x\})\le\tfrac1{24}-\eta$ whenever $3^x\varepsilon\ge18\cdot3^\delta$,
for $r$ close to $1$.  By Lemma~\ref{lem:eighteen}(c) take $c\ge2$ with
$K_c<18\cdot3^{\delta/2}$.  Let $x_0$ be the least integer with
$3^{x_0}\varepsilon\ge18\cdot3^\delta$; as $r\to1^+$, $x_0\to\infty$ and
$x_0\varepsilon\to0$.  For $x\ge x_0>c$, Lemma~\ref{lem:eighteen}(b) and its
monotonicity bound $\tau(\{x\})$ by
\[
  t(c)+\frac{35}{24}\cdot
  \frac{\ell(c)+\lambda_7(c)(\tfrac37)^{x_0}}{h_{x_0-2}(\tfrac15,\tfrac17,1)}
  \cdot\frac{r^{x_0-1}}{18\cdot3^\delta},
\]
which tends to $t(c)+\ell(c)/(18\cdot3^\delta)
=\tfrac1{24}-(\tfrac1{24}-t(c))(1-K_c/(18\cdot3^\delta))$, since
$h_k(\tfrac15,\tfrac17,1)\to(1-\tfrac15)^{-1}(1-\tfrac17)^{-1}=\tfrac{35}{24}$
and $r^{x_0}\to1$.  As $K_c/(18\cdot3^\delta)<3^{-\delta/2}$, any
$\eta<(\tfrac1{24}-t(c))(1-3^{-\delta/2})$ will do.

\emph{(ii) Near placements.}  For $r$ close to $1$, every $x\ge4$ with
$3^x\varepsilon\le18\cdot3^{-\delta}$ satisfies the hypotheses of
Lemma~\ref{lem:eighteen}(a), so
$t(x)-14\cdot3^{-x-3}\le\tau(\{x\})\le t(x)$.  Such $x$ have $x\le s$, and
$y^x(1+y)-y^2\ge2y^{x+1}-1\ge2r^{-s-1}-1\to1$.  By
Lemma~\ref{lem:eighteen}(c) there is $x_1$ with
$18\cdot3^{-x}s_x\le3^{\delta/2}$ for $x\ge x_1$, and then
$s_x\le3^{-\delta/2}/\varepsilon$; for $4\le x<x_1$,
$s_x\le3^{x_1}/8$ by Lemma~\ref{lem:threenode}(c).  Either way
$s_x+7\le1/(r\varepsilon)$ once $\varepsilon$ is small.

(a) Let $\delta\le\omega\le1-\delta$.  Then
$3^{x^*}\varepsilon=18\cdot3^{-\omega}\le18\cdot3^{-\delta}$ and
$x^*\to\infty$, so by (ii) $\tau(\{x^*\})$ is the tail of a full certificate
with zero set $\{1,x^*,z\}$, $z\ge x^*+3$, and
\begin{equation}\label{eq:oscworst}
  \tfrac5{36}(\tfrac35)^{x^*}R(x^*)\le\tfrac1{24}-\tau(\{x^*\})
  \le\tfrac5{36}(\tfrac35)^{x^*}R(x^*)+\tfrac{14}{27}3^{-x^*},
\end{equation}
where $R(x^*)\to1$ and $3^{-x^*}=5^{-x^*}(\tfrac35)^{x^*}$.  Every other
placement costs less for $r$ close to $1$.  A placement $x\ge x^*+1$ has
$3^x\varepsilon\ge54\cdot3^{-\omega}\ge18\cdot3^\delta$, so costs at most
$\tfrac1{24}-\eta$ by (i).  Placements $x\le3$ cost at most
$\tfrac1{48}$.  For $4\le x\le x^*-2$,
$\tfrac1{24}-\tau(\{x\})\ge\tfrac5{36}(\tfrac35)^x\cdot\tfrac9{20}
\ge\tfrac54\cdot\tfrac5{36}(\tfrac35)^{x^*}$, and for $x=x^*-1$,
$\tfrac1{24}-\tau(\{x\})\ge\tfrac53\cdot\tfrac5{36}(\tfrac35)^{x^*}R(x^*-1)$.
All of these exceed the right side of \eqref{eq:oscworst} for $r$ close to
$1$.  So $T=\tau(\{x^*\})$, \cite[Theorem~3.1]{coefficient-mass-attainment} and (ii) give the bounds
on the infimum, and by \eqref{eq:oscworst} and \eqref{eq:oscscale},
$576(\tfrac1{24}-T)\varepsilon^{-\alpha}=\Phi(\omega)\bigl(R(x^*)+O(5^{-x^*})\bigr)$,
which differs from $\Phi(\omega)\le\Phi(1)$ by $o(1)$ uniformly in $\omega$.

(b) \emph{Lower bound.}  For $r$ close to $1$ every placement $x$ has
$\tfrac1{24}-\tau(\{x\})$ at least $\eta$ if $3^x\varepsilon\ge18\cdot3^\delta$
(by (i)), at least $\tfrac1{48}$ if $x\le3$, at least
$\tfrac54\cdot\tfrac5{36}(\tfrac35)^s$ if $4\le x\le s-2$, and at least
$\tfrac5{36}(\tfrac35)^{s+\delta}R(x)$ if $s-2<x<s+\delta$, where
$R(x)\to1$ uniformly as $x>s-2\to\infty$.  So
$\tfrac1{24}-T=\inf_x(\tfrac1{24}-\tau(\{x\}))$ and \eqref{eq:oscscale} give
$\liminf_{r\to1^+}\Delta(r)\ge\Phi(0)(\tfrac35)^\delta$.

\emph{Upper bound.}  Let $\hat x$ be the largest integer with
$3^{\hat x}\varepsilon\le18\cdot3^{-\delta}$, so $\hat x>s-1-\delta$ and
$\hat x\to\infty$.  By (ii),
$\tfrac1{24}-T\le\tfrac1{24}-\tau(\{\hat x\})\le\tfrac5{36}(\tfrac35)^{\hat x}R(\hat x)
+\tfrac{14}{27}3^{-\hat x}$, where
$80(\tfrac35)^{\hat x}\varepsilon^{-\alpha}<\Phi(0)(\tfrac53)^{1+\delta}
=\Phi(1)(\tfrac53)^\delta$ by \eqref{eq:oscscale} and
$3^{-\hat x}\varepsilon^{-\alpha}<3^{1+\delta}\varepsilon^{1-\alpha}/18\to0$.
So $\limsup_{r\to1^+}\Delta(r)\le\Phi(1)(\tfrac53)^\delta$.

As $\delta$ is arbitrary, $\Phi(0)\le\liminf\Delta\le\limsup\Delta\le\Phi(1)$.
For $0<\omega_0<1$ the points $r_n=1+18\cdot3^{-n-\omega_0}$ have $s=n+\omega_0$,
$x^*=n$ and $\omega=\omega_0$, so $\Delta(r_n)\to\Phi(\omega_0)$ by (a).
Hence every point of $[\Phi(0),\Phi(1)]$ is a limit point, and the
$\liminf$ and $\limsup$ are the endpoints.
\end{proof}

At $r=\tfrac{5701}{5700}$, far below $\rho_1$, everything is exact:
$3^{10}\varepsilon=10.35\ldots$ and $3^{11}\varepsilon=31.07\ldots$, so
$x^*=10$ and $\omega=0.50\ldots$.  The hypotheses of
Lemma~\ref{lem:eighteen}(a) hold at $x=10$, and the witness of the full
certificate with zero set $\{1,10,1335\}$ has $s_1=0.4525\ldots$ and
$s_{1335}=-0.3405\ldots$, so by \eqref{eq:witness} $\tau(\{10\})$ is its
tail, which lies within $10^{-600}$ below $t(10)$; the placements $4\le x\le9$
cost at most $t(x)<\tau(\{10\})$, those $x\le3$ at most $\tfrac1{48}$, and
Lemma~\ref{lem:eighteen}(b) with $c=4$ at $x=11$ bounds every $x\ge11$ by
$0.03593<\tau(\{10\})=0.04085\ldots$.  So $T=\tau(\{10\})$ and
$\inf_Qb_2(QP_r)=24.47884\ldots$, with $\Delta(r)=26.70\ldots$ against
$\Phi(\omega)=26.97\ldots$; this is checked exactly.
There $t(11)>t(10)$, so the placement $11$ is beaten only through the node
$\tfrac1r$, while Lemma~\ref{lem:eighteen}(b) at $x=10$ beats $t(10)$ for no
$c$, as $3^{10}\varepsilon<18$; the test checks both as controls, and that
the hypothesis $s_x+7\le1/(r\varepsilon)$ of Lemma~\ref{lem:eighteen}(a)
fails at $x=11$.
Theorem~\ref{thm:familyosc} leaves out neighbourhoods of the jumps, where
$3^{x^*}\varepsilon$ is just below $18$ or $3^{x^*+1}\varepsilon$ just above
it.  To cross them we first compute every cost along the family.

\section{The cost of one placement}\label{sec:onezero}

The witness of
Lemma~\ref{lem:eighteen}(a) flips its entries far out.  Flipping entries
just below $x$ as well gives the cost of every placement along the family
in closed form, with its minimising zero set, at every $1<r<3$.  Fix
$1<r<3$, $y=\tfrac1r$ and an integer $x\ge2$, and let $U=U^{(x)}$,
$t=t(x)$, $s_1$, $s_x$ and $\sigma$ be as in Lemma~\ref{lem:threenode},
whose parts (a) and (b) hold for $x\ge2$; also $0<s_1\le\tfrac12$, by
part (c) for $x\ge4$, while $s_1=\tfrac7{24}$ and $\tfrac{859}{3408}$ at
$x=2,3$.  Let $V$ and $V'$ be the exponential sums on
$(\tfrac13,\tfrac15,\tfrac17)$ with $(V_0,V_1,V_x)=(0,1,0)$ and
$(V'_0,V'_1,V'_x)=(0,0,1)$.  With $A,B,C,D$ as in
Lemma~\ref{lem:threenode}, their weights are $-\tfrac{105}{2D}(B-C,C-A,A-B)$
and $\tfrac1D(3,-10,7)$, so
$V'_d=h_{d-2}(\tfrac13,\tfrac15,\tfrac17)/h_{x-2}(\tfrac13,\tfrac15,\tfrac17)$
by the proof of Lemma~\ref{lem:threenode}(a).  Put
\[
  H_d=y^d-U_d-yV_d-y^xV'_d,\qquad
  F=\sum_{d\ge1}\sigma_dy^d-t=s_1y+s_xy^x+\frac{y^x+y^{x+1}-y^2}{1-y}-t .
\]
Call the integers $d\ge2$, $d\ne x$, the \emph{positions}; unmarked sums
over $d$ run over them.  Order them as $\ldots,x+2,x+1,x-1,\ldots,2$, and for
a position $z$ let $\Pi(z)$ be the set of positions before $z$ in this
order: all $d>z$ if $z>x$, and all $d>x$ with all $z<d<x$ if $z<x$.

\begin{theorem}[One placement along the family]\label{thm:familyzero}
In this notation:
\begin{enumerate}
\item[(a)] $H$ is the exponential sum on $(y,\tfrac13,\tfrac15,\tfrac17)$
with weight $1$ on $y$ vanishing at $d=0,1,x$.  At every position $U_d$ and
$H_d$ have the sign of $d-x$ and $V_d$ the opposite sign, and
\[
  \sum_d|V_d|=s_1,\qquad\sum_d|H_d|=F>0 .
\]
\item[(b)] If $0\le\varphi_d\le1$ at the positions and
$2\sum_d\varphi_d|H_d|\ge F$, then $\tau(\{x\})\ge t-2\sum_d\varphi_d|U_d|$.
\item[(c)] Let $z$ be the position with
$\sum_{d\in\Pi(z)}|H_d|<\tfrac F2\le\sum_{d\in\Pi(z)}|H_d|+|H_z|$, and
$\varphi=(\tfrac F2-\sum_{d\in\Pi(z)}|H_d|)/|H_z|\in(0,1]$.  Then
\[
  \tau(\{x\})=t-2\sum_{d\in\Pi(z)}|U_d|-2\varphi|U_z|,
\]
the tail of the full certificate on $(y,\tfrac13,\tfrac15,\tfrac17)$ with
zero set $\{1,x,z\}$ if $z>x$ and $\{1,z,x\}$ if $z<x$.  If $\varphi<1$, it is
the only certificate achieving $\tau(\{x\})$.
\end{enumerate}
\end{theorem}

So the third zero is a half-mass point of $|H|$, flipped from the far end,
as $N$ is in \cite[Theorem~5.3(c)]{coefficient-mass-attainment}; it lies above $x$ exactly when
$\sum_{d>x}H_d\ge\tfrac F2$.  Parts (b) and (c) are the witness of
Lemma~\ref{lem:eighteen}(a) with the correction at the positions $1$ and
$x$ written through the kernel $H$.

\begin{proof}
Write $\mathbf a=(\tfrac13,\tfrac15,\tfrac17)$ and $a$ for one of its
nodes.  The $3\times3$ matrix $(a^d)$, $d\in\{0,1,x\}$, is nonsingular by the
zero bound, so for every integer $d$ there are unique $c_0,c_1,c_x$ with
$a^d=c_0+c_1a+c_xa^x$ at the three nodes.  Summing against the weights of
$U$, which satisfy $\sum_a\lambda_a=1$ and $\sum_a\lambda_aa=\sum_a\lambda_aa^x=0$,
gives $c_0=U_d$, and likewise $c_1=V_d$, $c_x=V'_d$:
\begin{equation}\label{eq:threebasis}
  a^d=U_d+V_da+V'_da^x\qquad\text{at }a=\tfrac13,\tfrac15,\tfrac17 .
\end{equation}

(a) So $H_d=y^d-\sum_a\mu_aa^d$ with $\mu_a$ the weight of $U+yV+y^xV'$ at
$a$, and $H_0=1-1=0$, $H_1=y-y=0$, $H_x=y^x-y^x=0$; also $\sum_a\mu_a=1$.
By the zero bound, $U$ is negative on $(1,x)$ and positive past $x$
(proof of Lemma~\ref{lem:threenode}(a)).  The sum $V$ is nonzero on three
nodes and vanishes at $0$ and $x$, so it has no other zeros and changes
sign at both; as $V_1=1$, it is positive on $(0,x)$ and negative past $x$.
The sum $H$ is nonzero on four nodes and vanishes at $0$, $1$ and $x$, so it
has no other zeros and changes sign at each; its weight on the largest node
$y>\tfrac13$ is $1$, so it is positive past $x$ and negative on $(1,x)$.  By
Lemma~\ref{lem:threenode}(b), $\sum_{d\ge1}\sigma_da^d=t$ at the three
nodes; summing against the weights of $V$, which add up to $V_0=0$, gives
$\sum_{d\ge1}\sigma_dV_d=0$, and against those of $H$,
$\sum_{d\ge1}\sigma_dH_d=\sum_{d\ge1}\sigma_dy^d-t\sum_a\mu_a=F$.  At
every position $\sigma_d$ is $-1$ below $x$ and $1$ above it, so
$\sigma_dV_d=-|V_d|$ and $\sigma_dH_d=|H_d|$; with $V_1=1$, $V_x=0$ and
$H_1=H_x=0$ the two sums read $s_1-\sum_d|V_d|=0$ and $\sum_d|H_d|=F$.

To see $F>0$, put $p(a)=(1-a)(\sum_{d\ge1}\sigma_da^d-t)$, the polynomial
\[
  p(a)=-t+(s_1+t)a-(s_1+1)a^2+(s_x+1)a^x+(1-s_x)a^{x+1}
\]
for $x\ge3$, the terms in $a^2$ combining for $x=2$.  It has at most
five terms, so by Descartes' rule at most four positive zeros counted with
multiplicity.  It vanishes at $\tfrac17,\tfrac15,\tfrac13$, while
$p(0)=-t<0<1=p(1)$, so the number of its zeros in $(0,1)$ is odd, hence
three: $\tfrac17,\tfrac15,\tfrac13$, all simple.  So $p>0$ on
$(\tfrac13,1]$, and $F=p(y)/(1-y)>0$.

(b) Put $\eta_d=-2\varphi_d\sgn(U_d)$ at the positions.  For $0\le s\le1$
let $\rho(s)$ agree with $\sigma+s\eta$ at the positions and have
$\rho_1=s_1-s\sum_d\eta_dV_d$ and $\rho_x=s_x-s\sum_d\eta_dV'_d$; the series
converge absolutely, as $V$ and $V'$ decay geometrically.  By
\eqref{eq:threebasis}, at each node $a$,
\[
  \sum_{d\ge1}\rho_da^d=t+s\sum_d\eta_d(a^d-V_da-V'_da^x)
  =t+s\sum_d\eta_dU_d=t-2s\sum_d\varphi_d|U_d|=:\theta(s),
\]
and at $y$, as $\sgn H_d=\sgn U_d$,
\[
  \sum_{d\ge1}\rho_dy^d-\theta(s)=F+s\sum_d\eta_d(y^d-U_d-V_dy-V'_dy^x)
  =F-2s\sum_d\varphi_d|H_d| .
\]
This is affine in $s$, positive at $s=0$ and at most $0$ at $s=1$, so it
vanishes at some $s\in(0,1]$.  There $\rho=\rho(s)$ satisfies
$\sum_{d\ge1}\rho_da^d=\theta$ at all four nodes, with
$\theta=\theta(s)\ge t-2\sum_d\varphi_d|U_d|$.  At a position
$\rho_d=\sigma_d(1-2s\varphi_d)\in[-1,1]$, and $\sgn(U_d)V_d=-|V_d|$ gives
$\rho_1=s_1-2s\sum_d\varphi_d|V_d|\in[-s_1,s_1]$ by (a).  So $|\rho_d|\le1$
for every $d\ne x$.  For a certificate $v$ for $\{x\}$ on the four nodes,
summing over them as in the proof of \cite[Lemma~5.1]{coefficient-mass-attainment} gives
$\theta=\sum_{d\ge1}\rho_dv_d=\sum_{d\ne x}\rho_dv_d\le\Tail(v)$, so
$\tau(\{x\})\ge\theta$.

(c) As $z$ runs through the positions in their order, the sums
$\sum_{\Pi(z)}|H_d|$ increase from $0$ (the $H_d$ are summable) by the
steps $|H_z|>0$ and end at $F-|H_2|$, so $z$ exists, is unique, and
$0<\varphi\le1$.  Take $\varphi_d=1$ on $\Pi(z)$, $\varphi_z=\varphi$ and
$\varphi_d=0$ otherwise.  Then $2\sum_d\varphi_d|H_d|=F$, so the proof of
(b) runs with $s=1$, and $\rho=\rho(1)$ has $\rho_d=-\sigma_d$ on $\Pi(z)$,
$\rho_z=\sigma_z(1-2\varphi)$ and $\rho_d=\sigma_d$ at the other positions:
if $z>x$, $\rho_d$ is $-1$ on $(1,x)$, $1$ on $(x,z)$ and $-1$ past $z$; if
$z<x$, it is $-1$ on $(1,z)$, $1$ on $(z,x)$ and $-1$ past $x$.  The full
certificate $v$ with zero set $\{1,x,z\}$ has $v_0=1$ and changes sign
exactly at its three zeros, so $\rho_d=\sgn(v_d)$ off its zero set, and
$\theta=\sum_{d\ge1}\rho_dv_d=\Tail(v)$.  Hence
$\Tail(v)=\theta\le\tau(\{x\})\le\Tail(v)$, and $\theta$ is the displayed
value.  If $\varphi<1$ and $v'$ achieves $\tau(\{x\})$, then
$\theta=\sum_{d\ne x}\rho_dv'_d\le\sum_{d\ne x}|v'_d|=\theta$ forces
$\rho_dv'_d=|v'_d|$ for $d\ne x$, so $v'_d=0$ where $|\rho_d|<1$: at $d=1$,
as $|\rho_1|\le s_1<1$, and at $d=z$, as $|1-2\varphi|<1$.  So $v'$ vanishes
on $\{1,x,z\}$ and is $v$.
\end{proof}

With \cite[Theorem~5.3]{coefficient-mass-attainment}, this makes $T$ along the family a finite
computation in closed form at every $1<r<3$: $Z^*=\{1,3\}$, so
$T=\max(\tau^*,\tau(\{x\}):x=2\text{ or }4\le x<N)$, each $\tau(\{x\})$
given by (c) through finitely many powers of $r$ and geometric sums.  The
costs found in \cite[Theorem~5.4]{coefficient-mass-attainment} and Theorems~\ref{thm:family}
and~\ref{thm:familylow} are instances: at $r=\tfrac54$ the placement $4$ has
$z=5$, zero set $\{1,4,5\}$ and cost $\tfrac{7627}{314024}$, and at
$r=\tfrac{11}{10}$ the placement $5$ has $z=3$, zero set $\{1,3,5\}$ and
cost $\tfrac{96935}{3398808}$.  For a fixed placement $x$ the zero set can
change with $r$ only where $\tfrac F2=\sum_{\Pi(z)}|H_d|+|H_z|$ for a
position $z$, where $\{1,x,z\}$ (or $\{1,z,x\}$) and the certificate with
the next position in the order tie; each side of this equation is a
combination of powers of $r$ and of geometric sums in $r$, so these are
the boundaries of the pieces in closed form.  In the pattern after
Theorem~\ref{thm:familylow}, the entry $s_z$ of the witness is
$1-2\varphi$ for $z>x$, and passing from $-1$ to $1$ is $\varphi$ passing
from $1$ to $0$.

\paragraph{The limit profile.}  Theorem~\ref{thm:familyzero} settles what
the costs tend to on the scale $3^x(r-1)\to\kappa$.  Below, $g_d$ is the
sum $3^{-d}-\tfrac{10}3\,5^{-d}+\tfrac73\,7^{-d}=\tfrac8{315}h_{d-2}(\tfrac13,\tfrac15,\tfrac17)$
of the proof of Lemma~\ref{lem:threenode}(a), positive for $d\ge2$;
$U^\infty_d=\tfrac72\,7^{-d}-\tfrac52\,5^{-d}$ is the full certificate on
$(\tfrac15,\tfrac17)$ with zero set $\{1\}$, negative for $d\ge2$, with
tail $\tfrac1{24}$ (proof of \cite[Proposition~5.5]{coefficient-mass-attainment}); and
$M_c=2\sum_{d\ge c}g_d=3\cdot3^{-c}-\tfrac{25}3\,5^{-c}+\tfrac{49}9\,7^{-c}$.

\begin{theorem}[The limit profile]\label{thm:familylimit}
Let $\tau(\{x\})$ be the costs at $P_r=(x-r)(x-3)(x-5)(x-7)$, $u=1$, let
$t(c)$, $\ell(c)$ and $K_c$ be as in Lemmas~\ref{lem:threenode}
and~\ref{lem:eighteen}, and put $k_2=\infty$ and
\[
  k_c=\frac{18}{1-18\cdot3^{-c}+30\cdot5^{-c}-14\cdot7^{-c}}\qquad(c\ge3).
\]
\begin{enumerate}
\item[(a)] $k_3=\tfrac{66150}{1957}=33.80\ldots$, the $k_c$ decrease
strictly, and $k_c\to18$.  Let $G(\kappa)=\tfrac1{24}$ for
$0\le\kappa\le18$ and $G(\kappa)=t(c)+\ell(c)/\kappa$ for
$k_c\le\kappa\le k_{c-1}$, $c\ge3$.  Then $G$ is well defined, continuous and
nonincreasing, $G(\kappa)<\tfrac1{24}$ for $\kappa>18$, and
\[
  G(\kappa)=\min\Bigl(\frac1{24},\ \inf_{c\ge2}\Bigl(t(c)+\frac{\ell(c)}\kappa\Bigr)\Bigr)
  \qquad(\kappa>0).
\]
\item[(b)] If $x\to\infty$ and $r\to1^+$ with $3^x(r-1)\to\kappa\in[0,\infty)$,
then $\tau(\{x\})\to G(\kappa)$.
\item[(c)] In (b), if $\kappa<18$, the zero set of
Theorem~\ref{thm:familyzero}(c) is eventually $\{1,x,z\}$ with $z>x$ and
$(r-1)z\to\log\frac{36}{18+\kappa}$; if $k_c<\kappa<k_{c-1}$, it is eventually
$\{1,c,x\}$.
\item[(d)] There are constants $C_2>C_1>0$ with
$C_1\eta^{1+\alpha}\le\tfrac1{24}-G(18(1+\eta))\le C_2\eta^{1+\alpha}$ for
$0<\eta\le1$, where $\alpha=\log(5/3)/\log3$.
\item[(e)] As $\eta\to0^+$,
$(\tfrac1{24}-G(18(1+\eta)))(18/\eta)^{1+\alpha}$ has no limit: its limit
points fill $[\tfrac54,f(\theta^*)]=[1.25,1.3814\ldots]$, where
$f(\theta)=\tfrac54(2\theta-1)\theta^{-1-\alpha}$ and
$\theta^*=\tfrac{1+\alpha}{2\alpha}$.
\end{enumerate}
\end{theorem}

So $\tau(\{x\})$ does tend to $\min_c(t(c)+\ell(c)/\kappa)$, the limit of
the bounds of Lemma~\ref{lem:eighteen}(b), once $\kappa>18$; the minimum is
at the $c\ge3$ with $k_c\le\kappa\le k_{c-1}$ (so $c=3$ for
$\kappa\ge33.80\ldots$, $c=4$ for $21.95\ldots\le\kappa\le33.80\ldots$),
never at $c=2$, and $c\to\infty$ as $\kappa\downarrow18$.  At
$r=\tfrac{101}{100}$ the placement $7$ has $3^7\varepsilon=21.87$, just
below $k_4$, and zero set $\{1,4,7\}$ rather than the limiting $\{1,5,7\}$.

\begin{proof}
\emph{The profile.}  Write $A=3^{-c}$, $B=5^{-c}$, $C=7^{-c}$ and $D$ as in
Lemma~\ref{lem:threenode}.  Expanding the closed forms of
Lemma~\ref{lem:threenode}(a),
\begin{equation}\label{eq:profileid}
  |U^\infty_c|=\ell(c)\,g_c,\qquad
  \frac1{k_c}=\frac1{18}-M_{c+1},\qquad
  t(c)+\ell(c)\Bigl(\frac1{18}-M_c\Bigr)=\frac1{24}-2\sum_{d\ge c}|U^\infty_d| :
\end{equation}
the first is $\tfrac12(5B-7C)=\tfrac{3(5B-7C)}{2D}\cdot\tfrac D3$, the second is
$M_{c+1}=A-\tfrac53B+\tfrac79C$, and in the third both sides minus
$\tfrac1{24}$ equal $(98C-75B)/12$, the left one after the factor
$D=3A-10B+7C$ cancels.  As $g_d>0$, $M_c$ decreases strictly to $0$, with
$M_c-M_{c+1}=2g_c$, and $M_3=\tfrac{19}{315}>\tfrac1{18}>M_4$; so the $k_c$,
$c\ge3$, are positive, decrease strictly and tend to $18$.  Define $L(w)=0$
for $w\le0$ and
\[
  L(w)=2\sum_{d>c}|U^\infty_d|+\ell(c)(w-M_{c+1})\qquad
  (M_{c+1}\le w\le M_c,\ c\ge3).
\]
By the first identity the pieces agree at their common ends, where both
give $2\sum_{d\ge c}|U^\infty_d|$, so $L$ is continuous on $(-\infty,M_3]$
(at $0$ as $\sum_{d>c}|U^\infty_d|\to0$), nondecreasing, and positive on
$(0,M_3]$.  For $\kappa>0$ put $w(\kappa)=\tfrac1{18}-\tfrac1\kappa<M_3$.
If $\kappa\le18$ then $L(w(\kappa))=0$.  If $k_c\le\kappa\le k_{c-1}$ then
$M_{c+1}\le w(\kappa)\le M_c$ by the second identity, and
$\tfrac1{24}-L(w(\kappa))$ is affine in $1/\kappa$ with slope $\ell(c)$ and, by
the third, equal to $t(c)+\ell(c)/\kappa$ at $1/\kappa=\tfrac1{18}-M_c$;
so it is $t(c)+\ell(c)/\kappa$.  Hence $G(\kappa)=\tfrac1{24}-L(w(\kappa))$
for every $\kappa>0$, which is well defined, continuous, nonincreasing, and
below $\tfrac1{24}$ for $\kappa>18$; and $k_3=\tfrac{66150}{1957}$.

\emph{Asymptotics.}  Let $x\to\infty$ and $r\to1^+$ with
$3^x\varepsilon\to\kappa$, $\varepsilon=r-1$; then $x\varepsilon\to0$ and
$y^x\to1$.  Use the notation of Theorem~\ref{thm:familyzero} at $x$, and
write $A,B,C,D$ at $x$ now.  By Lemma~\ref{lem:eighteen}(c)
$3^{-x}s_x\to\tfrac1{18}$, and $1-y=\varepsilon/r$, so
\begin{equation}\label{eq:Fscale}
  (1-y)F=(1-y)(s_1y+s_xy^x-t)+y^x+y^{x+1}-y^2\to1+\tfrac\kappa{18}.
\end{equation}
Summing the geometric part of $\sum_{d>x}H_d$,
\[
  E:=F-2\sum_{d>x}H_d
  =s_1y+(s_x+1)y^x-\frac{y^2}{1-y}-t+2\sum_{d>x}(U_d+yV_d+y^xV'_d),
\]
where $0<\sum_{d>x}U_d\le t$, and $V_d<0$ for $d>x$ with
$\sum_{d>x}|V_d|\le s_1$, by Theorem~\ref{thm:familyzero}(a), and
$0<\sum_{d>x}V'_d=(\tfrac32A-\tfrac52B+\tfrac76C)/D<1$ for $x\ge3$.  As
$y^2/(1-y)=1/(r\varepsilon)$, this gives $E\le s_x+4-1/(r\varepsilon)$, and,
if $\kappa>0$, $3^{-x}E\to\tfrac1{18}-\tfrac1\kappa$.  Below $x$, $V_d>0$,
$-1<U_d<0$ and $y^d<1$, so for $2\le d<x$
\begin{equation}\label{eq:Hbelow}
  y^xV'_d-2\le|H_d|\le y^xV'_d+1,\qquad
  V'_d=\frac{3g_d}D=\frac{3^xg_d}{1-\tfrac{10}3(\tfrac35)^x+\tfrac73(\tfrac37)^x},
\end{equation}
hence $3^{-x}\sum_{c<d<x}|H_d|\to\sum_{d>c}g_d=\tfrac12M_{c+1}$ and
$3^{-x}|H_c|\to g_c$ for fixed $c\ge2$.  Finally the weights of $U^{(x)}$
tend to $(0,-\tfrac52,\tfrac72)$, so $\sum_{d\ge2}|U^{(x)}_d-U^\infty_d|\to0$.

(b) \emph{Lower bound.}  If $\kappa<18$, then
$r\varepsilon(s_x+4)\to\kappa/18<1$, so eventually $E<0$, that is,
$2\sum_{d>x}H_d>F$; Theorem~\ref{thm:familyzero}(b) with $\varphi_d=1$ for
$d>x$ gives $\tau(\{x\})\ge t(x)-2\sum_{d>x}|U^{(x)}_d|\to\tfrac1{24}$.  Now
let $\kappa>0$, $w=w(\kappa)$, and $\max(w,0)<w'\le M_3$; take $c\ge3$ with
$M_{c+1}<w'\le M_c$ and $\varphi=(w'-M_{c+1})/(2g_c)\in(0,1]$, and put
$\varphi_d=1$ for $d>x$ and for $c<d<x$, and $\varphi_c=\varphi$.  Then
$3^{-x}(2\sum_d\varphi_d|H_d|-2\sum_{d>x}H_d)\to M_{c+1}+2\varphi g_c=w'$,
which exceeds $w=\lim3^{-x}E$, so eventually $2\sum_d\varphi_d|H_d|\ge F$,
and Theorem~\ref{thm:familyzero}(b) gives
\[
  \tau(\{x\})\ge t(x)-2\sum_{d>x}|U_d|-2\sum_{c<d<x}|U_d|-2\varphi|U_c|
  \to\frac1{24}-2\sum_{d>c}|U^\infty_d|-2\varphi|U^\infty_c|=\frac1{24}-L(w'),
\]
by the first identity of \eqref{eq:profileid}.  Letting $w'\downarrow\max(w,0)$,
$\liminf\tau(\{x\})\ge\tfrac1{24}-L(\max(w,0))=G(\kappa)$ in every case.

\emph{Upper bound.}  $\tau(\{x\})\le t(x)\to\tfrac1{24}$, and for $c\ge2$
and $x>c$ the bound of Lemma~\ref{lem:eighteen}(b) is
$t(c)+\tfrac{35}{24}(\ell(c)+\lambda_7(c)(\tfrac37)^x)r^{x-1}/(3^x\varepsilon\,h_{x-2}(\tfrac15,\tfrac17,1))$,
which tends to $t(c)+\ell(c)/\kappa$ if $\kappa>0$, as
$h_{x-2}(\tfrac15,\tfrac17,1)\to\tfrac{35}{24}$.  So $\limsup\tau(\{x\})$ is at
most $\min(\tfrac1{24},\inf_c(t(c)+\ell(c)/\kappa))\le G(\kappa)$, the last as
$G(\kappa)$ is $\tfrac1{24}$ or one of the $t(c)+\ell(c)/\kappa$.  So
$\tau(\{x\})\to G(\kappa)$, and, as every $\kappa>0$ is such a limit (take
$r=1+\kappa3^{-x}$), $G(\kappa)$ is that minimum; this proves (b) and the
rest of (a).

(c) If $\kappa<18$, eventually $E<0$, that is, $\sum_{d>x}H_d>\tfrac F2$, so
$z>x$ and $\sum_{d>z}H_d<\tfrac F2\le\sum_{d\ge z}H_d$.  For $z>x$,
$(1-y)\sum_{d\ge z}H_d=y^z-(1-y)\sum_{d\ge z}(U_d+yV_d+y^xV'_d)=y^z+O(\varepsilon)$
uniformly, by the bounds above, so by \eqref{eq:Fscale}
$y^{z+1}+o(1)<\tfrac12(1+\tfrac\kappa{18})\le y^z+o(1)$; hence
$y^z\to\tfrac{18+\kappa}{36}$, and $z\log r\to\log\tfrac{36}{18+\kappa}$.
If $k_c<\kappa<k_{c-1}$, then $M_{c+1}<w<M_c$, while $3^{-x}E\to w$,
$3^{-x}\cdot2\sum_{c<d<x}|H_d|\to M_{c+1}$ and
$3^{-x}\cdot2\sum_{c\le d<x}|H_d|\to M_c$; as $E=F-2\sum_{d>x}H_d$, eventually
$\sum_{\Pi(c)}|H_d|<\tfrac F2\le\sum_{\Pi(c)}|H_d|+|H_c|$, so $z=c$.

(d) Let $0<\eta\le1$, $\kappa=18(1+\eta)$ and $w=w(\kappa)=\eta/(18(1+\eta))$.
\emph{Upper.}  Let $c\ge3$ with $M_{c+1}<w\le M_c$.  On $[0,w]$ the function
$L$ is affine on pieces with slopes $\ell(c')$, $c'\ge c$, and
$\ell(c')\le5(\tfrac35)^{c'}$ ($\ell(3)=\tfrac{81}{142}$, and
$\ell(c')\le\tfrac{15B}{2D}\le\tfrac{75}{17}(\tfrac35)^{c'}$ for $c'\ge4$,
as $D\ge\tfrac{17}{10}A$ there by the proof of
Lemma~\ref{lem:threenode}(c)).  Also
$w>M_{c+1}\ge3^{-c}(1-\tfrac53(\tfrac35)^c)\ge\tfrac58\,3^{-c}$, so
$(\tfrac35)^c=(3^{-c})^\alpha\le(\tfrac85w)^\alpha$, and
$\tfrac1{24}-G(\kappa)=L(w)\le5(\tfrac85w)^\alpha w\le8(\eta/18)^{1+\alpha}$.
\emph{Lower.}  By (a), $\tfrac1{24}-G(\kappa)\ge\tfrac1{24}-t(c)-\ell(c)/\kappa
=(\tfrac1{24}-t(c))(1-K_c/\kappa)$ for every $c\ge2$.  For $c\ge4$, by the
proof of Lemma~\ref{lem:eighteen}(c), $K_c=18/(1-q)$ with
$q=(45AB-84AC+35BC)/(5B-7C)\le(45A+35C)/(5-7(\tfrac57)^c)\le14.6A\le\tfrac{14.6}{81}$,
so $K_c\le18(1+18\cdot3^{-c})$.  Take the least $c\ge4$ with
$18\cdot3^{-c}\le\eta/2$; then $3^{-c}>\eta/108$, $1-K_c/\kappa\ge
1-(1+\tfrac\eta2)/(1+\eta)\ge\tfrac\eta4$, and
$\tfrac1{24}-t(c)\ge\tfrac1{16}(\tfrac35)^c>\tfrac1{16}(\eta/108)^\alpha$ by
Lemma~\ref{lem:threenode}(c), so $\tfrac1{24}-G(\kappa)\ge\tfrac\eta{64}(\eta/108)^\alpha$.

(e) Let $w\to0^+$, let $c$ be the piece with $M_{c+1}<w\le M_c$, so
$c\to\infty$, and put $\theta=3^cw$; as $3^cM_{c+1}\to1$ and
$3^cM_c\to3$, the limit points of $\theta$ lie in $[1,3]$.  Here
$2\sum_{d>c}|U^\infty_d|=\tfrac54\,5^{-c}-\tfrac76\,7^{-c}$ and
$\ell(c)=\tfrac52(\tfrac35)^c(1+o(1))$, so
$L(w)=5^{-c}(\tfrac54+\tfrac52(\theta-1)+o(1))$, while
$w^{1+\alpha}=\theta^{1+\alpha}5^{-c}$ as $3^{1+\alpha}=5$.  Hence
$L(w)/w^{1+\alpha}-f(\theta)\to0$.  Along $w=\theta_03^{-c}$, $c\to\infty$,
with $\theta_0\in(1,3)$ fixed, $\theta\to\theta_0$; and $f(1)=f(3)=\tfrac54$,
while on $[1,3]$ the function $f$ has its maximum at $\theta^*$, where
$f'$ vanishes.  So the limit points of $L(w)/w^{1+\alpha}$ fill
$[\tfrac54,f(\theta^*)]$.  With $w=\eta/(18(1+\eta))$ the quantity in (e)
is $L(w)w^{-1-\alpha}(1+\eta)^{-1-\alpha}$, with the same limit points.
\end{proof}

\section{Across the jumps}\label{sec:jumps}

Near a jump the placements $x^*$ and $x^*+1$
compete on the scale $(\tfrac35)^{x^*}$, and the profile $G$ decides
between them.

\begin{theorem}[Across the jumps]\label{thm:familyjump}
In the notation of Theorems~\ref{thm:familyosc} and~\ref{thm:familylimit},
extend $\Phi(\omega)=80\cdot18^{-\alpha}(\tfrac53)^\omega$ to all real $\omega$, and put
\[
  \Gamma(r)=576\,\varepsilon^{-\alpha}\Bigl(\frac1{24}-G\bigl(3^{x^*+1}\varepsilon\bigr)\Bigr),
  \qquad\nu(r)=(1-\omega)\,\varepsilon^{-\alpha/(1+\alpha)} .
\]
\begin{enumerate}
\item[(a)] For $r$ close to $1$ every worst placement is $\{x^*\}$ or
$\{x^*+1\}$, and, with no restriction on $\omega$,
\[
  \Delta(r)-\min\bigl(\Phi(\omega),\,\Phi(\omega-1)+\Gamma(r)\bigr)\to0
  \qquad(r\to1^+).
\]
\item[(b)] There are constants $a_2>a_1>0$ such that, for $r$ close to $1$,
$\{x^*\}$ is the only worst placement if $\nu(r)\ge a_2$, and $\{x^*+1\}$ is
the only one if $\nu(r)\le a_1$.
\item[(c)] As $r\to1^+$: $\Delta(r)-\Phi(\omega)\to0$ if $\nu(r)\to\infty$,
in particular uniformly on $0<\omega\le1-\delta$ for each $\delta>0$; and
$\Delta(r)-\Phi(\omega-1)\to0$ if $\nu(r)\to0$.
\end{enumerate}
\end{theorem}

Here $\Phi(\omega-1)=\tfrac35\Phi(\omega)$, and $\Gamma(r)$ is of order
$\nu(r)^{1+\alpha}$ when $1-\omega\le\tfrac12$ (proof of (b)), with a ratio
that oscillates as $1-\omega\to0$ by Theorem~\ref{thm:familylimit}(e).  So as $r$
decreases through a jump, $\Delta$ does not jump from $\Phi(1)$ to $\Phi(0)$:
it passes from $\Phi(\omega)$ to $\Phi(\omega-1)$ while $1-\omega$ runs
through a window of order $(r-1)^{\alpha/(1+\alpha)}$,
$\alpha/(1+\alpha)=\log(5/3)/\log5=0.317\ldots$, which shrinks to the jump,
the worst placement switching from $x^*$ to $x^*+1$ inside it; the costs
there are those of Theorem~\ref{thm:familyzero}(c), with zero sets
$\{1,x,z\}$ or $\{1,z,x\}$.

\begin{proof}
Let $r\to1^+$; then $x^*\to\infty$, $s\varepsilon\to0$ and
$\varepsilon^{-\alpha}(\tfrac35)^{x^*}=\Phi(\omega)/80\le\Phi(1)/80$.  Write
$\beta_x=(\tfrac35)^x$ and $\gamma_x=(\tfrac37)^x$, and let $o(\beta)$ stand
for a quantity whose ratio to $\beta_{x^*}$ tends to $0$.  For
$x\in\{x^*,x^*+1\}$ put $\kappa=3^x\varepsilon\in[6,54]$ and
$w=\tfrac1{18}-\tfrac1\kappa$; so $x\varepsilon=o(\beta)$, $3^{-x}=o(\beta)$,
$\gamma_x=o(\beta)$, and $R(x)\to1$.  We show
\begin{align}
  \tfrac1{24}-\tau(\{x\})&\ge\tfrac5{36}\beta_x+\min\bigl(\tfrac1{24}-G(\kappa),\beta_x\bigr)-o(\beta),
  \label{eq:jumplow}\\
  \tfrac1{24}-\tau(\{x\})&\le\tfrac5{36}\beta_x+\bigl(\tfrac1{24}-G(\kappa)\bigr)+o(\beta)
  \quad\text{if }\tfrac1{24}-G(\kappa)\le\beta_x .\label{eq:jumpup}
\end{align}

\emph{Proof of \eqref{eq:jumplow}.}  As $\tau(\{x\})\le t(x)=\tfrac1{24}-\tfrac5{36}\beta_xR(x)$,
it holds when $\tfrac1{24}-G(\kappa)=o(\beta)$, in particular when
$\kappa\le18$.  Let $\kappa>18$ and $c\ge3$ with $k_c\le\kappa\le k_{c-1}$,
so $G(\kappa)=t(c)+\ell(c)/\kappa$.  If $c\ge x$, then
$0<w\le\tfrac1{18}-\tfrac1{k_{x-1}}=M_x$ and, the slopes of $L$ being at
most $1$ (proof of Theorem~\ref{thm:familylimit}(d)),
$\tfrac1{24}-G(\kappa)=L(w)\le M_x\le4\cdot3^{-x}=o(\beta)$.  Let $c<x$.  The
proof of Lemma~\ref{lem:eighteen}(b) gives
$\tau(\{x\})\le t(c)+J\cdot3^xU^{(c)}_x/\kappa$ with
$J=\tfrac{35}{24}r^{x-1}/h_{x-2}(\tfrac15,\tfrac17,1)$, where
$3^xU^{(c)}_x=\ell(c)+\lambda_5(c)\beta_x+\lambda_7(c)\gamma_x>0$,
$\lambda_5(c)$ being the weight of $U^{(c)}$ on $\tfrac15$.  Here $J\ge1$ and
$J-1=O(x\varepsilon+x5^{-x})=o(\beta)$, as
$\tfrac{35}{24}-h_k(\tfrac15,\tfrac17,1)=\sum_{i+j>k}5^{-i}7^{-j}\le(k+2)5^{-k}$;
$\lambda_5(c)\le-\tfrac52$ by Lemma~\ref{lem:threenode}(a), as
$10\cdot5^{-c}\ge14\cdot7^{-c}$; and
$\ell(c)\le1$, $0<\lambda_7(c)\le7$ for $c\ge3$.  So
$\tau(\{x\})\le G(\kappa)-\tfrac5{2\kappa}\beta_x+o(\beta)$.  Fix
$\zeta\in(0,1)$.  If $\kappa\le18(1+\zeta)$, then
$\tfrac5{2\kappa}\ge\tfrac5{36}(1-\zeta)$; otherwise
$\tfrac1{24}-G(\kappa)\ge\tfrac1{24}-G(18(1+\zeta))>0$, a fixed quantity,
which exceeds $\tfrac{41}{36}\beta_x$ for $r$ close to $1$.  Either way
$\tfrac1{24}-\tau(\{x\})\ge\tfrac5{36}\beta_x+\min(\tfrac1{24}-G(\kappa),\beta_x)
-\tfrac5{36}\zeta\beta_x-o(\beta)$, and $\zeta$ is arbitrary.

\emph{Proof of \eqref{eq:jumpup}.}  Now $L(\max(w,0))=\tfrac1{24}-G(\kappa)\le\beta_x$,
so by Theorem~\ref{thm:familylimit}(d) $\max(w,0)\to0$.  By the bound
$E\le s_x+4-1/(r\varepsilon)$ of the proof of
Theorem~\ref{thm:familylimit}, with
$18\cdot3^{-x}s_x=(1-36A+90B-56C)/(1-\tfrac{10}3\beta_x+\tfrac73\gamma_x)\le1+4\beta_x$
and $3^{-x}/(r\varepsilon)=1/(r\kappa)\ge(1-\varepsilon)/\kappa$, we get
$3^{-x}E\le w+\tfrac14\beta_x$ for $r$ close to $1$.  Let
$w'=\max(w,0)+\tfrac12\beta_x$, take $c'\ge3$ with $M_{c'+1}<w'\le M_{c'}$
and $\varphi=(w'-M_{c'+1})/(2g_{c'})$, and put $\varphi_d=1$ for $d>x$ and
for $c'<d<x$, and $\varphi_{c'}=\varphi$.  As $w'\to0$, $c'\to\infty$; as
$M_{c'}\ge w'\ge\tfrac12\beta_x$ and $M_{c'}\le4\cdot3^{-c'}$,
$3^{c'}\le8/\beta_x$, so $c'<x$.  By \eqref{eq:Hbelow},
$|H_d|\ge3^xg_d(1-\xi)-2$ for $2\le d<x$, with
$\xi=1-y^x/(1+\tfrac73\gamma_x)=o(\beta)$, so
\[
  2\sum_d\varphi_d|H_d|-2\sum_{d>x}H_d
  \ge3^x(1-\xi)(w'-M_x)-4x\ge3^x(w'-\tfrac14\beta_x)\ge E
\]
for $r$ close to $1$, that is, $2\sum_d\varphi_d|H_d|\ge F$.  By
Theorem~\ref{thm:familyzero}(b),
$\tau(\{x\})\ge t(x)-2\sum_{d>c'}|U_d|-2\varphi|U_{c'}|$, with
$U=U^{(x)}$.  Its weights differ from $(0,-\tfrac52,\tfrac72)$ by at most
$15\beta_x$ each ($D\ge\tfrac{17}{10}A$), so
$|U_d|\le|U^\infty_d|+45\beta_x3^{-d}$ and the loss is at most
$L(w')+135\beta_x3^{-c'}=L(w')+o(\beta)$.  Finally
$L(w')-L(\max(w,0))\le\tfrac12\beta_x\cdot5(\tfrac35)^{c'}=o(\beta)$, the
slopes of $L$ on $[\max(w,0),w']$ being $\ell(c'')\le5(\tfrac35)^{c''}$ with
$c''\ge c'$.  So
$\tau(\{x\})\ge t(x)-(\tfrac1{24}-G(\kappa))-o(\beta)$, which is
\eqref{eq:jumpup}.

\emph{The other placements.}  Placements $x\le3$ cost at most
$\tfrac1{48}$, and $4\le x<x^*$ at most $t(x)$, with
$\tfrac1{24}-t(x)=\tfrac5{36}\beta_xR(x)$, at least
$\tfrac5{36}\cdot\tfrac53R(x^*-1)\beta_{x^*}$ for $x=x^*-1$ and
$\tfrac5{36}\cdot\tfrac{25}9\cdot\tfrac9{20}\beta_{x^*}$ below; both exceed
$\tfrac5{36}\cdot\tfrac98\beta_{x^*}$ for $r$ close to $1$.  Placements
$x\ge x^*+2$ have $3^x\varepsilon\ge54$ and cost at most
$\tfrac1{24}-\eta$ for a fixed $\eta>0$ by step (i) of the proof of
Theorem~\ref{thm:familyosc}.

(a) At $x^*$, $\kappa<18$ and $G(\kappa)=\tfrac1{24}$, so by
\eqref{eq:jumplow} and \eqref{eq:jumpup}
$\tfrac1{24}-\tau(\{x^*\})=\tfrac5{36}\beta_{x^*}+o(\beta)$; every placement
$x\ne x^*,x^*+1$ falls short of $\tau(\{x^*\})$ by a multiple of
$\beta_{x^*}$, so every worst placement is $\{x^*\}$ or $\{x^*+1\}$ and
$\tfrac1{24}-T=\min(\tfrac1{24}-\tau(\{x^*\}),\tfrac1{24}-\tau(\{x^*+1\}))$.
At $x^*+1$ put $g=\tfrac1{24}-G(3^{x^*+1}\varepsilon)$.  If
$g\le\beta_{x^*+1}$, \eqref{eq:jumplow} and \eqref{eq:jumpup} give
$\tfrac1{24}-\tau(\{x^*+1\})=\tfrac5{36}\beta_{x^*+1}+g+o(\beta)$.  If
$g>\beta_{x^*+1}$, \eqref{eq:jumplow} gives
$\tfrac1{24}-\tau(\{x^*+1\})\ge\tfrac{41}{36}\cdot\tfrac35\beta_{x^*}-o(\beta)$,
more than $\tfrac5{36}\beta_{x^*}$ by a multiple of $\beta_{x^*}$, as is
$\tfrac5{36}\beta_{x^*+1}+g$.  Either way
\[
  \tfrac1{24}-T=\min\bigl(\tfrac5{36}\beta_{x^*},\ \tfrac5{36}\beta_{x^*+1}+g\bigr)+o(\beta).
\]
Multiply by $576\varepsilon^{-\alpha}$: by \eqref{eq:oscscale}
$80\beta_{x^*}\varepsilon^{-\alpha}=\Phi(\omega)$ and
$80\beta_{x^*+1}\varepsilon^{-\alpha}=\Phi(\omega-1)$, and
$576\varepsilon^{-\alpha}g=\Gamma(r)$.  As in the proof of
Theorem~\ref{thm:familyosc}, $\Delta(r)-576(\tfrac1{24}-T)\varepsilon^{-\alpha}\to0$,
which gives (a).

(b) Let $v=1-\omega$ and $\eta=3^v-1$, so $3^{x^*+1}\varepsilon=18(1+\eta)$.
If $v\le\tfrac12$, then $v\le\eta\le2v\le1$, and
Theorem~\ref{thm:familylimit}(d) gives
$576C_1\nu^{1+\alpha}\le\Gamma(r)\le2304C_2\nu^{1+\alpha}$, as
$\varepsilon^{-\alpha}v^{1+\alpha}=\nu^{1+\alpha}$; if $v>\tfrac12$, then
$g\ge\tfrac1{24}-G(18\sqrt3)>0$ and $\Gamma(r)\to\infty$.  Choose $a_2$ with
$576C_1a_2^{1+\alpha}\ge\Phi(1)+1$ and $a_1$ with
$2304C_2a_1^{1+\alpha}\le\tfrac1{10}\Phi(0)$.  If $\nu\ge a_2$ (and $r$ is close
to $1$), then $g>\beta_{x^*+1}$ or $\Gamma\ge\Phi(1)+1$, and in both cases
the estimates above give $\tfrac1{24}-\tau(\{x^*+1\})$ at least
$\tfrac1{24}-\tau(\{x^*\})$ plus a fixed multiple of
$\varepsilon^\alpha$, since $576\varepsilon^{-\alpha}\beta_{x^*+1}=\tfrac{36}5\Phi(\omega-1)\ge90>\Phi(1)+1$;
so $\{x^*\}$ is the only worst placement.  If $\nu\le a_1$, then
$g\le\tfrac1{5760}\Phi(0)\varepsilon^\alpha<\beta_{x^*+1}$, and
$576\varepsilon^{-\alpha}(\tfrac1{24}-\tau(\{x^*+1\}))\le\Phi(\omega-1)+\tfrac1{10}\Phi(0)+o(1)
\le\Phi(\omega)-\tfrac3{10}\Phi(0)+o(1)$, while
$576\varepsilon^{-\alpha}(\tfrac1{24}-\tau(\{x^*\}))=\Phi(\omega)+o(1)$; so
$\{x^*+1\}$ is the only worst placement.

(c) If $\nu\to\infty$ then $\Gamma\to\infty$ by the bounds in (b), and if
$\nu\to0$ then $\Gamma\to0$; (a) gives the two limits, as
$\Phi(\omega-1)<\Phi(\omega)$.  On $0<\omega\le1-\delta$,
$\nu\ge\delta\varepsilon^{-\alpha/(1+\alpha)}\to\infty$.
\end{proof}

Two exact examples show the switch.  At $r=\tfrac{101}{100}$,
$3^6\varepsilon=7.29$ and $3^7\varepsilon=21.87$, so $x^*=6$, $\omega=0.82\ldots$
and $\nu(r)=0.76\ldots$.  By Theorem~\ref{thm:familyzero}(c) the placement
$6$ costs the tail $0.035595\ldots$ of $\{1,6,41\}$, and the placement $7$
the tail $0.035944\ldots$ of $\{1,4,7\}$; the placements $4,5$ cost at most
$t(x)<t(6)$, those $x\le3$ at most $\tfrac1{48}$, and
Lemma~\ref{lem:eighteen}(b) with $c=3$ at $x=8$ bounds every $x\ge8$ by
$0.02761$.  So $T=\tau(\{7\})$ and $\inf_Qb_2(QP_r)=27.8209\ldots$: the
worst placement is $x^*+1$, as the floating-point search after
Theorem~\ref{thm:familylow} suggested.  At
$r=\tfrac{9842}{9841}$, $3^{10}\varepsilon=6.0003\ldots$ and
$3^{11}\varepsilon=18.0009\ldots$, so $x^*=10$ and $1-\omega=4.6\ldots\cdot10^{-5}$,
deep in the window.  The placement $11$ costs the tail $0.0411505\ldots$ of
$\{1,7,11\}$, more than $t(10)=0.0408515\ldots$ and than $t(x)$ for
$4\le x\le9$, while Lemma~\ref{lem:eighteen}(b) with $c=3$ at $x=12$ bounds
every $x\ge12$ by $0.0288$.  So $T=\tau(\{11\})$,
$\inf_Qb_2(QP_r)=24.3009\ldots$, and $\Delta(r)=21.63\ldots$, near
$\Phi(\omega-1)=20.86\ldots$ rather than $\Phi(\omega)=34.77\ldots$.  Both are
checked exactly.

Theorem~\ref{thm:familyosc}(a) holds on an explicit range.

\begin{proposition}[An explicit range]\label{prop:familyrdelta}
Let $0<\delta\le\tfrac12$.  If $1<r\le1+(\delta/10)^4$ and
$\delta\le\omega\le1-\delta$, in the notation of
Theorem~\ref{thm:familyosc}, then $\{x^*\}$ is the only worst placement,
$T=\tau(\{x^*\})$ is the tail of the full certificate with zero set
$\{1,x^*,z\}$ for some $z\ge x^*+3$, and
$1/t(x^*)\le\inf_Qb_2(QP_r)\le1/(t(x^*)-14\cdot3^{-x^*-3})$.
\end{proposition}

\begin{proof}
Write $x=x^*$, $\kappa=3^x\varepsilon=18\cdot3^{-\omega}$ and
$\beta=(\tfrac35)^x$, and let $\varepsilon\le(\delta/10)^4\le20^{-4}$.  We
use $0.4649<\alpha<0.465$, so $1/\alpha<2.151$, $\log_3\tfrac73>0.77$ and
$\tfrac12\log_35>0.73$.  Each bound below of a power of $\varepsilon$ by a
multiple of $\delta$ is used at $\varepsilon=(\delta/10)^4$, where its two
sides have a ratio increasing in $\delta$ (a positive power of $\delta$,
times a logarithm for the bound on $x\varepsilon$); so it suffices to check
it at $\delta=\tfrac12$.  As
$3^{x+1}\varepsilon\ge18$, $3^x\ge6/\varepsilon\ge960000$, so $x\ge13$ and
$3^{-x}\le\varepsilon/6$; hence $\beta=(3^{-x})^\alpha\le(\varepsilon/6)^\alpha
\le\delta/250$, $(\tfrac37)^x\le(\varepsilon/6)^{0.77}$ and
$5^{-x/2}\le(\varepsilon/6)^{0.73}$.  As $3^x<18/\varepsilon$,
$x\varepsilon<\varepsilon\log_3(18/\varepsilon)\le\delta/2000$, the middle
term increasing in $\varepsilon$ and at most $\delta/2000$ at
$\varepsilon=(\delta/10)^4$.

\emph{The placement $x$.}  With $y^x\ge1-x\varepsilon$,
$y^x(1+y)-y^2\ge(1-x\varepsilon)(2-\varepsilon)-1>\varepsilon>\tfrac{1-y}{24}$.
By the proof of Lemma~\ref{lem:eighteen}(c),
$18\cdot3^{-x}s_x\le(1+90\cdot5^{-x})/(1-\tfrac{10}3\beta)\le1+4\beta$ for
$x\ge13$, so $s_x\le3^{-\delta}(1+4\beta)/\varepsilon$ as
$\kappa\le18\cdot3^{-\delta}$; with $3^{-\delta}\le1-\tfrac56\delta$ on
$[0,\tfrac12]$ (convexity, $3^{-1/2}<\tfrac7{12}$) and
$4\beta+9\varepsilon<\tfrac56\delta$, this gives
$(s_x+7)r\varepsilon\le(1-\tfrac56\delta+4\beta+7\varepsilon)(1+\varepsilon)<1$.
So Lemma~\ref{lem:eighteen}(a) applies: $\tau(\{x\})$ is the tail of
$\{1,x,z\}$ with $z\ge x+3$, and $t(x)-14\cdot3^{-x-3}\le\tau(\{x\})\le t(x)$.
For $x\ge11$ the closed form of $R(x)$ in Lemma~\ref{lem:eighteen}(c) has
numerator in $[1-\tfrac75(\tfrac57)^{11}-9\cdot3^{-11}-7^{-10},1+\tfrac{84}5\,3^{-11}]$
and denominator in $[1-\tfrac{10}3(\tfrac35)^{11},1+\tfrac73(\tfrac37)^{11}]$,
so $0.96\le R(x)\le1.02$, and
$\tfrac1{24}-\tau(\{x\})\le\tfrac5{36}\cdot1.02\,\beta+14\cdot3^{-x-3}\le0.143\beta$.

\emph{Placements below $x$.}  Those $x'\le3$ cost at most $\tfrac1{48}$, and
$4\le x'<x$ at most $t(x')$, with
$\tfrac1{24}-t(x-1)=\tfrac5{36}\cdot\tfrac53R(x-1)\beta\ge0.22\beta$ and
$\tfrac1{24}-t(x')\ge\tfrac1{16}(\tfrac35)^{x'}\ge\tfrac1{16}\cdot\tfrac{25}9\beta>0.17\beta$
for $x'\le x-2$ (Lemma~\ref{lem:threenode}(c)).  All are below
$\tau(\{x\})$.

\emph{Placements above $x$.}  Let $c$ be the least integer $c\ge4$ with
$18\cdot3^{-c}\le\delta/2$, so $3^c<108/\delta<3^x$, and put
$\kappa'=3^{x+1}\varepsilon=18\cdot3^{1-\omega}\ge18\cdot3^\delta\ge18(1+\tfrac{21}{20}\delta)$.
By Lemma~\ref{lem:eighteen}(b) and its monotonicity every $x'>x$ costs at
most $t(c)+\ell(c)\Pi/\kappa'$, where
\[
  \Pi=\Bigl(1+\frac{\lambda_7(c)}{\ell(c)}\bigl(\tfrac37\bigr)^{x+1}\Bigr)\,
  r^x\,\frac{35}{24\,h_{x-1}(\tfrac15,\tfrac17,1)} .
\]
Here, with $A,B,C$ at $c$, $\lambda_7(c)/\ell(c)=\tfrac{7(3A-5B)}{3(5B-7C)}\le\tfrac{7A}{5B-7C}\le2.3(\tfrac53)^c
<2.3(108/\delta)^\alpha$, and $(\tfrac37)^x\le(\varepsilon/6)^{0.77}$,
so the first factor is at most $1+\delta/100$; $r^x\le e^{x\varepsilon}\le1+\delta/1000$;
and $h_{x-1}(\tfrac15,\tfrac17,1)\ge\tfrac{35}{24}(1-5^{-x/2})(1-7^{-x/2})$,
the terms with both exponents at most $\tfrac{x-1}2$, so the last factor is
at most $1+3\cdot5^{-x/2}\le1+\delta/100$.  So $\Pi\le1+\delta/30$.  By the
proof of Theorem~\ref{thm:familylimit}(d), $K_c\le18(1+18\cdot3^{-c})\le18(1+\tfrac\delta2)$,
so
\[
  \frac1{24}-t(c)-\frac{\ell(c)\Pi}{\kappa'}
  =\Bigl(\frac1{24}-t(c)\Bigr)\Bigl(1-\frac{K_c\Pi}{\kappa'}\Bigr)
  \ge\frac1{16}\Bigl(\frac35\Bigr)^c\Bigl(1-\frac{(1+\tfrac\delta2)(1+\tfrac\delta{30})}{1+\tfrac{21}{20}\delta}\Bigr)
  \ge0.02\,\delta\Bigl(\frac35\Bigr)^c .
\]
This exceeds $0.143\beta\ge\tfrac1{24}-\tau(\{x\})$ once
$(\tfrac35)^{x-c}\le\delta/7.2$, that is, $x-c\ge\log_{5/3}(7.2/\delta)$,
which holds as $c<\log_3(108/\delta)$, $x\ge\log_3(6/\varepsilon)$ and
$(\delta/10)^4\le\tfrac\delta{18}(\tfrac\delta{7.2})^{1/\alpha}$ for
$0<\delta\le\tfrac12$ ($1/\alpha=2.150\ldots$).  So every placement above
$x$ costs less than $\tau(\{x\})$.

Hence $\{x\}$ is the only worst placement, $T=\tau(\{x\})$, and
\cite[Theorem~3.1]{coefficient-mass-attainment} gives the bounds on the infimum.
\end{proof}

So the escaping placement stays extremal at some roots below $2$, here down
to $r_c$, and the failure below $r_c$ starts at the placement $\{4\}$ just
above $Z^*$, where the conclusion of \cite[Lemma~4.1]{coefficient-mass-attainment} fails.  The cutoff cannot be
lowered uniformly, however.

\section{Open questions}\label{sec:familyopen}

Along $(r,3,5,7)$ the infimum is known on $[1.0745\ldots,3)$ in closed form
(Theorems~\ref{thm:family} and~\ref{thm:familylow}); below, every cost is
explicit at each $r$ (Theorem~\ref{thm:familyzero}), infinitely many pieces
accumulate at $r=1$ (Corollary~\ref{cor:familyone}), and near $r=1$ the
infimum is described uniformly: $(\inf_Qb_2(QP_r)-24)(r-1)^{-\alpha}$ is
$\min(\Phi(\omega),\Phi(\omega-1)+\Gamma(r))+o(1)$, with $\Gamma$ given
by the limit profile $G$ (Theorems~\ref{thm:familylimit}
and~\ref{thm:familyjump}), and the conclusions of
Theorem~\ref{thm:familyosc}(a) hold for $r\le1+(\delta/10)^4$
(Proposition~\ref{prop:familyrdelta}).  Open:
\begin{itemize}
\item the order of the pieces on $(1,1.0745\ldots)$, that is, whether the
worst placement is nonincreasing in $r$ and whether each boundary equation
after Theorem~\ref{thm:familyzero} has a single root in $r$;
\item explicit constants $a_1,a_2$ and a rate of convergence in
Theorem~\ref{thm:familyjump};
\item since $(\tfrac1{24}-G(18(1+\eta)))\eta^{-1-\alpha}$ oscillates
(Theorem~\ref{thm:familylimit}(e)), the window itself is log-periodic:
whether $\Delta$ is monotone in $r$ across each window, and where exactly
the worst placement switches, in the scale $\nu$;
\item beyond this family, an a priori bound on the zero set minimising one
$\tau(\{x\})$ at general roots.  Here Theorem~\ref{thm:familyzero} gives it
exactly, using that the minimiser at the three fast nodes is $\{1,x\}$ with
a witness whose corrections at $1$ and $x$ keep their signs.
\end{itemize}

\paragraph{Data availability.}
No datasets were generated or analyzed in this study.  The exact
computations, certificates and reproduction instructions are available in
the accompanying versioned source archive.

\paragraph{Computer assistance.}
Five computations enter the arguments: exact symbolic identities and sign
checks in the proofs of Theorems~\ref{thm:family} and
\ref{thm:familylow}, and the three rational examples following
Theorems~\ref{thm:familyosc} and~\ref{thm:familyjump}.  The symbolic
identities are verified over $\mathbb Q(r)$, interval signs are certified by
Sturm sequences over $\mathbb Q$, and the examples use exact rational
linear algebra and finite comparisons.  Appendix~\ref{sec:familydata}
records the certificate data needed to inspect the two theorem proofs; the
versioned source archive contains the complete exact computations.  No
floating-point calculation is used as a proof premise.

\paragraph{Use of AI tools.}
Large language model assistants (Claude, Anthropic) were used in drafting
and revising the text, the proofs and the verification code.  The author has checked the mathematics and is responsible
for the content.

\paragraph{Funding and competing interests.}
The author received no specific funding for this work and has no competing
interests to declare.

\appendix

\section{Data for the family \texorpdfstring{$(r,3,5,7)$}{(r,3,5,7)}}\label{sec:familydata}

This appendix lists the solutions of the linear systems in the proof of
Theorem~\ref{thm:family}, and the data of Theorem~\ref{thm:familylow}, as
rational functions of $r$, written with
\[
  p(r)=(3-r)(5-r)(7-r),\qquad q(r)=3466r^2+7455r+11025 ,
\]
so that the tails $1/\beta(r)$,
$t_2(r)=(4267r^2+5070r-5871)/\bigl(48(r-1)q(r)\bigr)$ and $t_5(r)$ there can be checked
by hand or by any computer algebra system.  Let $1<r<3$, so that
$\tfrac1r$ is the largest of the nodes
$\tfrac1r,\tfrac13,\tfrac15,\tfrac17$; then $p(r)$ and $q(r)$ are positive.  Each sum below has the form
$v_d=\lambda_1r^{-d}+\lambda_33^{-d}+\lambda_55^{-d}+\lambda_77^{-d}$, and
past its largest prescribed zero $z$ its values sum to
\[
  \sum_{d>z}v_d=\frac{\lambda_1}{r^{z}(r-1)}+\frac{\lambda_3}{2\cdot3^{z}}
  +\frac{\lambda_5}{4\cdot5^{z}}+\frac{\lambda_7}{6\cdot7^{z}} .
\]
Each weight vector is verified by substituting it into the defining
conditions; the values listed are then direct evaluations.  The signs are
those of the zero bound: a full certificate is positive at $0$ and changes
sign exactly at its zeros, and $H$ vanishes exactly at $0,1,3$, with
$H_d>0$ for $d>3$.  The escaping minimiser $V$, with zero set $Z^*=\{1,3\}$,
is displayed after \cite[Theorem~4.3]{coefficient-mass-attainment}.

\paragraph{The certificate with zero set $\{1,3,4\}$.}  Its weights are
\[
  \begin{aligned}
  \lambda_1&=-\frac{15r^4}{p(r)(15r+71)},&
  \lambda_3&=\frac{81(r+12)}{8(3-r)(15r+71)},\\
  \lambda_5&=-\frac{625(r+10)}{4(5-r)(15r+71)},&
  \lambda_7&=\frac{2401(r+8)}{8(7-r)(15r+71)}.
  \end{aligned}
\]
and $v_0=1$, $v_1=v_3=v_4=0$,
\[
  v_2=-\frac1{15r+71},\qquad
  \sum_{d\ge5}v_d=-\frac{r+14}{48(r-1)(15r+71)},
\]
with $v_2<0$ and $v_d<0$ for $d\ge5$, so
\[
  \Tail(v)=|v_2|-\sum_{d\ge5}v_d=\frac{49r-34}{48(r-1)(15r+71)}=\frac1{\beta(r)}.
\]

\paragraph{The certificate with zero set $\{1,2,5\}$.}  Its weights are
\begin{gather*}
  \lambda_1=-\frac{3466r^5}{p(r)q(r)},\qquad
  \lambda_3=\frac{243(109r^2+420r+1225)}{8(3-r)q(r)},\\
  \lambda_5=-\frac{3125(79r^2+210r+441)}{4(5-r)q(r)},\qquad
  \lambda_7=\frac{16807(49r^2+120r+225)}{8(7-r)q(r)},
\end{gather*}
and $v_0=1$, $v_1=v_2=v_5=0$,
\[
  v_3=\frac{71r+105}{q(r)},\qquad
  v_4=\frac{15r+71}{q(r)},\qquad
  \sum_{d\ge6}v_d=-\frac{139r^2+750r+2577}{48(r-1)q(r)},
\]
with $v_3,v_4>0$ and $v_d<0$ for $d\ge6$, so
\[
  \Tail(v)=v_3+v_4-\sum_{d\ge6}v_d
  =\frac{4267r^2+5070r-5871}{48(r-1)q(r)}=t_2(r).
\]

\paragraph{The certificate with zero set $\{1,3,5\}$.}  With
$q_5(r)=960r^2+5041r+7455$, its weights are
\begin{gather*}
  \lambda_1=-\frac{960r^5}{p(r)q_5(r)},\qquad
  \lambda_3=\frac{729(r+5)(r+7)}{2(3-r)q_5(r)},\\
  \lambda_5=-\frac{15625(r+3)(r+7)}{2(5-r)q_5(r)},\qquad
  \lambda_7=\frac{16807(r+3)(r+5)}{(7-r)q_5(r)},
\end{gather*}
and $v_0=1$, $v_1=v_3=v_5=0$,
\[
  v_2=-\frac{71r+105}{q_5(r)},\qquad
  v_4=\frac{r+15}{q_5(r)},\qquad
  \sum_{d\ge6}v_d=-\frac{7r^2+98r+375}{24(r-1)q_5(r)},
\]
with $v_2<0$, $v_4>0$ and $v_d<0$ for $d\ge6$, so
\[
  \Tail(v)=-v_2+v_4-\sum_{d\ge6}v_d
  =\frac{5(347r^2+250r-501)}{24(r-1)q_5(r)}=t_5(r).
\]

\paragraph{The sum $H$ of \cite[Theorem~5.3]{coefficient-mass-attainment}.}  With weight $1$ on
$\tfrac1r$, $H_0=0$ and $H|_{Z^*}=0$, its weights are
\[
  \begin{aligned}
  \lambda_1&=1,&
  \lambda_3&=-\frac{27(5-r)(7-r)(12r+35)}{568r^3},\\
  \lambda_5&=\frac{125(3-r)(7-r)(10r+21)}{284r^3},&
  \lambda_7&=-\frac{343(3-r)(5-r)(8r+15)}{568r^3}.
  \end{aligned}
\]
and $H_0=H_1=H_3=0$,
\begin{gather*}
  H_2=-\frac{p(r)}{71r^3},\qquad
  H_4=\frac{p(r)(15r+71)}{7455r^4},\qquad
  H_5=\frac{p(r)(960r^2+5041r+7455)}{782775r^5},\\
  \sum_{d\ge4}H_d=\frac{p(r)(14r+57)}{3408r^3(r-1)},\qquad
  \sum_{d\ge6}H_d=\frac{p(r)(10890r^3+50099r^2+80656r+119280)}{12524400r^5(r-1)} .
\end{gather*}
So $\Tail(H)=-H_2+\sum_{d\ge4}H_d$ and $\sum_{d\le5}|H_d|=-H_2+H_4+H_5$,
which give the identity for $2\sum_{d\le5}|H_d|-\Tail(H)$ in the proof of
Theorem~\ref{thm:family}.  For the identity with $d\le6$ there,
\[
  H_6=\frac{p(r)(44535r^3+246086r^2+529305r+782775)}{82191375r^6},\qquad
  \sum_{d\ge7}H_d=\frac{p(r)\,e_7(r)}{1315062000r^6(r-1)},
\]
with $e_7(r)=430890r^4+2035579r^3+3937376r^2+8468880r+12524400$, and
$\sum_{d\le6}|H_d|=-H_2+H_4+H_5+H_6$.

\paragraph{The certificates of Theorem~\ref{thm:familylow}.}  With
\begin{align*}
  Q_5&=77r^2+480r+1733,\qquad Q_6=9359r^3+60705r^2+246086r+363930,\\
  Q_7&=421939r^4+2798055r^3+12013156r^2+25839030r+38212650,\\
  Q_8&=8457547r^5+56817765r^4+251780638r^3+630690690r^2\\
  &\qquad+1356549075r+2006164125,
\end{align*}
the tails of the full certificates with zero sets $\{1,4,z\}$ are
\[
  \theta_5=\frac{n_5}{96(r-1)Q_5},\qquad
  \theta_6=\frac{n_6}{48(r-1)Q_6},\qquad
  \theta_7=\frac{n_7}{48(r-1)Q_7},\qquad
  \theta_8=\frac{n_8}{96(r-1)Q_8},
\]
where
\begin{align*}
  n_5&=769r^2+3374r-3989,\qquad n_6=49550r^3+244468r^2+148970r-433629,\\
  n_7&=2308315r^4+12063338r^3+17059705r^2+15420090r-46429509,\\
  n_8&=94332044r^5+508724104r^4+933286436r^3+1788992289r^2\\
  &\qquad+1607184810r-4915604589 .
\end{align*}
For $\{1,4,5\}$ the weights are
\begin{gather*}
  \lambda_1=-\frac{77r^5}{p(r)Q_5},\qquad
  \lambda_3=\frac{243(r^2+12r+109)}{16(3-r)Q_5},\\
  \lambda_5=-\frac{3125(r^2+10r+79)}{8(5-r)Q_5},\qquad
  \lambda_7=\frac{16807(r^2+8r+49)}{16(7-r)Q_5},
\end{gather*}
and $v_0=1$, $v_1=v_4=v_5=0$,
\[
  v_2=-\frac{15r+71}{2Q_5},\qquad v_3=-\frac{r+15}{2Q_5},\qquad
  \sum_{d\ge6}v_d=-\frac{r^2+14r+139}{96(r-1)Q_5},
\]
all negative, so $\Tail(v)=-v_2-v_3-\sum_{d\ge6}v_d=\theta_5(r)$.  Its witness
in \eqref{eq:witness} has
\begin{gather*}
  s_1=\frac{769r^3+10766r^2+42091r-52276}{96(r-1)Q_5},\qquad
  s_4=\frac{211252r^3+380157r^2+254562r-991257}{96(r-1)Q_5},\\
  s_5=-\frac{142079r^3+137846r^2-345779r-55456}{32(r-1)Q_5},
\end{gather*}
and $s_d=-1$ for $d=2,3$ and $d\ge6$.  The other witnesses have larger
coefficients and are only computed.  The polynomials of the breakpoints are
\begin{align*}
  \gamma_1&=218644r^3+418845r^2+374850r-1157625,\\
  \gamma_2&=203860r^3+341469r^2+134274r-824889,\\
  \gamma'&=172587443623r^6+203696380262r^5-245023502843r^4-16113960832r^3\\
  &\qquad-40364204160r^2-86819140800r-128394504000,\\
  \gamma_3&=31863386060r^7+527299949544r^6+2216976743064r^5+5148085143336r^4\\
  &\qquad+7473795889089r^3+2333772047565r^2-9236198868099r-16544067678495,\\
  \gamma_4&=1638534863r^5+1910488822r^4-2452391683r^3-384420992r^2\\
  &\qquad-826848960r-1222804800,\\
  \gamma_5&=15476503r^4+17470982r^3-26164523r^2-7874752r-11645760,\\
  \gamma_6&=144543r^3+150742r^2-305683r-110912,\qquad
  \gamma_7=3989r^2+3570r-11025,\\
  \gamma_+&=111700r^3-50307r^2-97470r-109209,
\end{align*}
each with exactly one zero in $(1,\tfrac32)$, where it changes sign from
$-$ to $+$.  Besides the ties after Theorem~\ref{thm:familylow},
\begin{gather*}
  \theta_5-\theta_6=-\frac{(r^2+15r+154)\gamma_6}{32(r-1)Q_5Q_6},\qquad
  \theta_6-\theta_7=-\frac{(15r^3+225r^2+2310r+9359)\gamma_5}{16(r-1)Q_6Q_7},\\
  \theta_7-\theta_8=-\frac{7(22r^4+330r^3+3388r^2+17745r+60277)\gamma_4}{32(r-1)Q_7Q_8} .
\end{gather*}

\begin{thebibliography}{9}

\bibitem{coefficient-mass-attainment}
B.~Pham,
\newblock Sharpness of the coefficient order-statistic bound,
\newblock preprint, 2026, \url{https://github.com/bangyen/coefficient-mass}.

\end{thebibliography}

\end{document}
